% !TEX program = lualatex
% ---------------------------------------------------------------------------------------------
%  Readback note 39: Links of pairs of terms: integral lifts and the discriminant gain
%
%  Build:  python3 tools/rb.py note <slug>
%
%  One unit of a Lean artifact, rewritten as ordinary mathematics. The prose is hand-written. Lean names,
%  line ranges and counts come from data.tex; quoted Lean statements are copied from the artifact by the build.
%
%  Rules of the text: standard terminology; every symbol is explained where it first appears and listed again
%  in the notation at the end; no letter has two meanings; a proof is complete or says what it takes for granted;
%  Lean names appear only in the LEAN lines and in the concordance.
%
%  The signature rule (the build enforces it): a function, relation or predicate that the artifact defines is
%  written only through its notation macro, which the build makes from the spec (`notation`) with one slot for
%  each explicit argument of the Lean definition. No argument is fixed and dropped, also not "for brevity" and
%  not for one section. Every argument and every hypothesis of a stated declaration stands inside its
%  statement, or in a Setting that the statement cites by number; only implicit and instance arguments (the
%  ambient objects) may stand in the conventions.
% ---------------------------------------------------------------------------------------------
\documentclass[11pt,a4paper]{article}
\ifdefined\pdfvariable\pdfvariable suppressoptionalinfo \numexpr 1+2+4+8\relax\fi
\usepackage[a4paper,hmargin=35mm,top=22mm,bottom=27mm,footskip=12mm,nohead]{geometry}

\usepackage{fontspec}
\directlua{
  luaotfload.add_fallback("gfallback", {"STIXTwoMath-Regular.otf:mode=node;", "DejaVuSans.ttf:mode=node;"})
  luaotfload.add_fallback("gmonofallback", {"DejaVuSansMono.ttf:mode=node;", "STIXTwoMath-Regular.otf:mode=node;"})
}
\setmainfont{NewCM10}[Extension=.otf,UprightFont=*-Book,ItalicFont=*-BookItalic,
  BoldFont=*-Bold,BoldItalicFont=*-BoldItalic,SmallCapsFont=*-Book,SmallCapsFeatures={Letters=SmallCaps},
  RawFeature={fallback=gfallback}]
\setsansfont{NewCMSans10}[Extension=.otf,UprightFont=*-Book,ItalicFont=*-BookOblique,
  BoldFont=*-Bold,BoldItalicFont=*-BoldOblique,RawFeature={fallback=gfallback}]
\setmonofont{JetBrainsMono}[Extension=.otf,UprightFont=*-Regular,ItalicFont=*-Italic,
  BoldFont=*-SemiBold,BoldItalicFont=*-SemiBoldItalic,Scale=MatchLowercase,
  RawFeature={-calt;-liga;fallback=gmonofallback}]
\usepackage{amsmath}
\usepackage{amsthm}
\usepackage[math-style=TeX]{unicode-math}
\setmathfont{NewCMMath-Book.otf}
\emergencystretch=2.5em\relax

\usepackage{xcolor}
\usepackage{tikz}
\usetikzlibrary{calc,arrows.meta,positioning,backgrounds}
\usepackage[protrusion=true,expansion=true]{microtype}
\usepackage{enumitem}
\usepackage{needspace}
\usepackage{booktabs}
\usepackage{array}
\usepackage{longtable}
\usepackage{fancyhdr}
\usepackage{lastpage}
\usepackage{multicol}
\usepackage[hidelinks,pdfusetitle]{hyperref}

\definecolor{ink}{HTML}{0F172A}
\definecolor{inksoft}{HTML}{475569}
\definecolor{inkfaint}{HTML}{94A3B8}
\definecolor{rulegrey}{HTML}{CBD5E1}
\definecolor{arith}{HTML}{C2410C}
\definecolor{est}{HTML}{7C3AED}
\definecolor{ana}{HTML}{0891B2}
\definecolor{ok}{HTML}{047857}
\definecolor{leancol}{HTML}{3730A3}

\title{Links of pairs of terms: integral lifts and the discriminant gain}
\author{grh-compact}
\date{}

% ------------------------------------------------------------------------------------- macros
\newcommand{\leanfont}{\ttfamily\fontsize{8.6}{10.5}\selectfont\color{leancol}}
\newcommand{\lean}[1]{\csname LN@#1\endcsname}
\newcommand{\fillin}[1]{{\color{arith}\itshape[#1]}}  % a place of the template that is not written yet; the build reports every one that is left
\newcommand{\leanquote}[1]{\csname LQ@#1\endcsname}
% generated by rb.py note; do not edit
\newcommand{\ConcDecls}{9}
\newcommand{\Pin}{72fead10}
\newcommand{\PinText}{that of commit {\leanfont 72fead10}}
\newcommand{\ArtifactName}{grh-compact}
\expandafter\newcommand\csname LN@absnormbordereddet\endcsname{\href{https://github.com/lsernau/grh-lean/blob/72fead102f663d5c7a407964a13dc0cf9d977207/grh-compact/src/New/Filling/LinkPairing.lean\#L25-L50}{{\leanfont Filling.\allowbreak{}HasLevel.\allowbreak{}abs\_norm\_bordered\_det}}}
\expandafter\newcommand\csname LN@existsliftslinkDiscSqle\endcsname{\href{https://github.com/lsernau/grh-lean/blob/72fead102f663d5c7a407964a13dc0cf9d977207/grh-compact/src/New/Filling/LinkPairing.lean\#L55-L122}{{\leanfont Filling.\allowbreak{}exists\_lifts\_linkDiscSq\_le}}}
\expandafter\newcommand\csname LN@existsisIntMatrix\endcsname{\href{https://github.com/lsernau/grh-lean/blob/72fead102f663d5c7a407964a13dc0cf9d977207/grh-compact/src/New/Filling/LinkPairing.lean\#L130-L138}{{\leanfont Filling.\allowbreak{}HasIntegralLifts.\allowbreak{}exists\_isIntMatrix}}}
\expandafter\newcommand\csname LN@existsorthogonalsplit\endcsname{\href{https://github.com/lsernau/grh-lean/blob/72fead102f663d5c7a407964a13dc0cf9d977207/grh-compact/src/New/Filling/LinkPairing.lean\#L151-L162}{{\leanfont Filling.\allowbreak{}HasIntegralLifts.\allowbreak{}exists\_orthogonal\_split}}}
\expandafter\newcommand\csname LN@existsmemlinktermDisceq\endcsname{\href{https://github.com/lsernau/grh-lean/blob/72fead102f663d5c7a407964a13dc0cf9d977207/grh-compact/src/New/Filling/LinkPairing.lean\#L167-L179}{{\leanfont Filling.\allowbreak{}HasLevel.\allowbreak{}exists\_mem\_link\_termDisc\_eq}}}
\expandafter\newcommand\csname LN@onelexRowsIndexandsaturationIndex\endcsname{\href{https://github.com/lsernau/grh-lean/blob/72fead102f663d5c7a407964a13dc0cf9d977207/grh-compact/src/New/Filling/LinkPairing.lean\#L181-L184}{{\leanfont Filling.\allowbreak{}one\_le\_xRowsIndex\_and\_saturationIndex}}}
\expandafter\newcommand\csname LN@minorNormsqlelinkDiscSqmul\endcsname{\href{https://github.com/lsernau/grh-lean/blob/72fead102f663d5c7a407964a13dc0cf9d977207/grh-compact/src/New/Filling/LinkPairing.lean\#L193-L232}{{\leanfont Filling.\allowbreak{}minorNorm\_sq\_le\_linkDiscSq\_mul}}}
\expandafter\newcommand\csname LN@existslinksplitleclassGauge\endcsname{\href{https://github.com/lsernau/grh-lean/blob/72fead102f663d5c7a407964a13dc0cf9d977207/grh-compact/src/New/Filling/LinkPairing.lean\#L239-L278}{{\leanfont Filling.\allowbreak{}exists\_link\_split\_le\_classGauge}}}
\expandafter\newcommand\csname LQ@existslinksplitleclassGauge\endcsname{\leanquoteline{0}{lemma exists\_link\_split\_le\_classGauge γ' : (0 < S.latticeDisc c (commonLink γ γ') ∧ 1 ≤}\leanquoteline{2}{S.frameValueIndex c (S.termPosition c c' γ) 1 (commonLink γ γ')) ∧ ∃ᵉ (i) (ρ) (\_ : Fintype ρ) (m :}\leanquoteline{2}{Matrix \_ ρ E) (y : Matrix \_ (Fin i) E) (v) (z : Matrix ρ \_ (𝔸f E)), let L}}
\expandafter\newcommand\csname LN@pairInvariantmulle\endcsname{\href{https://github.com/lsernau/grh-lean/blob/72fead102f663d5c7a407964a13dc0cf9d977207/grh-compact/src/New/Filling/LinkPairing.lean\#L284-L315}{{\leanfont Filling.\allowbreak{}pairInvariant\_mul\_le}}}
\expandafter\newcommand\csname LQ@pairInvariantmulle\endcsname{\leanquoteline{0}{lemma pairInvariant\_mul\_le (h' : γ' ∈ S.termsOver c s c') (hd : S.HasLevel j γ') : ∃ i, 𝔰.k ≤}\leanquoteline{2}{finrank E (commonLink γ γ') + i + (j - 𝔰.ϱ) ∧ (i = 0 → link γ = link γ') ∧}\leanquoteline{2}{S.pairInvariant c c' γ γ' * S.pairInvariant c c' γ' γ * (if 𝔰.ϱ < j ∧ i = 0 then x else 1) \textasciicircum{} 2 ≤}\leanquoteline{2}{(2 / S.g) \textasciicircum{} (2 * i * numInfPlaces E) * (S.classGauge j c c' γ * S.classGauge j c c' γ')}}
\newcommand{\ConcordanceRows}{\crow{absnormbordereddet}{\href{https://github.com/lsernau/grh-lean/blob/72fead102f663d5c7a407964a13dc0cf9d977207/grh-compact/src/New/Filling/LinkPairing.lean\#L25-L50}{{\leanfont Filling.\allowbreak{}HasLevel.\allowbreak{}abs\_norm\_bordered\_det}}}{LinkPairing.lean}{25--50}
\crow{existsliftslinkDiscSqle}{\href{https://github.com/lsernau/grh-lean/blob/72fead102f663d5c7a407964a13dc0cf9d977207/grh-compact/src/New/Filling/LinkPairing.lean\#L55-L122}{{\leanfont Filling.\allowbreak{}exists\_lifts\_linkDiscSq\_le}}}{LinkPairing.lean}{55--122}
\crow{existsisIntMatrix}{\href{https://github.com/lsernau/grh-lean/blob/72fead102f663d5c7a407964a13dc0cf9d977207/grh-compact/src/New/Filling/LinkPairing.lean\#L130-L138}{{\leanfont Filling.\allowbreak{}HasIntegralLifts.\allowbreak{}exists\_isIntMatrix}}}{LinkPairing.lean}{130--138}
\crow{existsorthogonalsplit}{\href{https://github.com/lsernau/grh-lean/blob/72fead102f663d5c7a407964a13dc0cf9d977207/grh-compact/src/New/Filling/LinkPairing.lean\#L151-L162}{{\leanfont Filling.\allowbreak{}HasIntegralLifts.\allowbreak{}exists\_orthogonal\_split}}}{LinkPairing.lean}{151--162}
\crow{existsmemlinktermDisceq}{\href{https://github.com/lsernau/grh-lean/blob/72fead102f663d5c7a407964a13dc0cf9d977207/grh-compact/src/New/Filling/LinkPairing.lean\#L167-L179}{{\leanfont Filling.\allowbreak{}HasLevel.\allowbreak{}exists\_mem\_link\_termDisc\_eq}}}{LinkPairing.lean}{167--179}
\crow{onelexRowsIndexandsaturationIndex}{\href{https://github.com/lsernau/grh-lean/blob/72fead102f663d5c7a407964a13dc0cf9d977207/grh-compact/src/New/Filling/LinkPairing.lean\#L181-L184}{{\leanfont Filling.\allowbreak{}one\_le\_xRowsIndex\_and\_saturationIndex}}}{LinkPairing.lean}{181--184}
\crow{minorNormsqlelinkDiscSqmul}{\href{https://github.com/lsernau/grh-lean/blob/72fead102f663d5c7a407964a13dc0cf9d977207/grh-compact/src/New/Filling/LinkPairing.lean\#L193-L232}{{\leanfont Filling.\allowbreak{}minorNorm\_sq\_le\_linkDiscSq\_mul}}}{LinkPairing.lean}{193--232}
\crow{existslinksplitleclassGauge}{\href{https://github.com/lsernau/grh-lean/blob/72fead102f663d5c7a407964a13dc0cf9d977207/grh-compact/src/New/Filling/LinkPairing.lean\#L239-L278}{{\leanfont Filling.\allowbreak{}exists\_link\_split\_le\_classGauge}}}{LinkPairing.lean}{239--278}
\crow{pairInvariantmulle}{\href{https://github.com/lsernau/grh-lean/blob/72fead102f663d5c7a407964a13dc0cf9d977207/grh-compact/src/New/Filling/LinkPairing.lean\#L284-L315}{{\leanfont Filling.\allowbreak{}pairInvariant\_mul\_le}}}{LinkPairing.lean}{284--315}}
% the notations of the defined functions: one slot for each explicit argument of the Lean declaration
\newcommand{\UnitMembers}{497}
\newcommand{\UnitInNote}{9}
\newcommand{\UnitProofs}{179}
\newcommand{\UnitSteps}{1371}
\newcommand{\UnitAllSteps}{1479}


\linespread{1.08}                       % 11 pt on 14.7 pt
\setlength{\parindent}{0pt}
\setlength{\parskip}{5.5pt plus 1.5pt minus 0.5pt}
\setlength{\abovedisplayskip}{7.5pt plus 2pt minus 3pt}
\setlength{\belowdisplayskip}{7.5pt plus 2pt minus 3pt}
\setlength{\abovedisplayshortskip}{3pt plus 2pt}
\setlength{\belowdisplayshortskip}{6pt plus 2pt minus 2pt}
\emergencystretch=1.5em
\clubpenalty=10000 \widowpenalty=10000 \displaywidowpenalty=10000
\relpenalty=7000 \binoppenalty=9000
\raggedbottom
\pagestyle{fancy}\fancyhf{}\renewcommand{\headrulewidth}{0pt}
\fancyfoot[C]{\sffamily\fontsize{7.6}{9}\selectfont\color{inkfaint}\thepage\,/\,\pageref*{LastPage}}
\setlist{topsep=2pt,partopsep=0pt,parsep=0pt,itemsep=4pt,leftmargin=2.1em,labelsep=0.55em}

\newcommand{\Q}{\mathbb{Q}}
\newcommand{\Z}{\mathbb{Z}}
\newcommand{\F}{\mathbb{F}}
\newcommand{\len}{\operatorname{len}}
\newcommand{\cok}{\operatorname{cok}}
\newcommand{\End}{\operatorname{End}}
\newcommand{\Hom}{\operatorname{Hom}}
\newcommand{\Gal}{\operatorname{Gal}}
\newcommand{\diag}{\operatorname{diag}}
\newcommand{\kw}[1]{\emph{#1}}
\newcommand{\m}[1]{\mbox{$#1$}}        % a short formula that must not be broken across lines
\newcommand{\head}[1]{\par\addvspace{17pt plus 4pt minus 2pt}\Needspace{6\baselineskip}%
  \noindent{\fontsize{13.2}{16}\selectfont\bfseries\color{ink}#1}\par\nopagebreak\vspace{-1pt}}
\newcommand{\sect}[2]{\refstepcounter{section}\label{#2}\par\addvspace{19pt plus 4pt minus 2pt}\Needspace{7\baselineskip}%
  \noindent\llap{{\fontsize{13.2}{16}\selectfont\bfseries\color{inkfaint}\thesection}\hspace{9pt}}%
  {\fontsize{13.2}{16}\selectfont\bfseries\color{ink}#1}\par\nopagebreak\vspace{-1pt}}
\newcommand{\subhead}[1]{\par\addvspace{9pt plus 2pt}\Needspace{4\baselineskip}\noindent\textbf{#1}\enspace\ignorespaces}
% where a statement lives in the Lean archive
\newcommand{\leantag}{{\sffamily\bfseries\fontsize{6.3}{8}\selectfont\color{leancol}\addfontfeatures{LetterSpace=7}LEAN}}
\newcommand{\sep}{\unskip\nobreak\hspace{5pt}{\color{inkfaint}\textperiodcentered}\hspace{5pt plus 2pt}\ignorespaces}
\newcommand{\leanline}[1]{\par\nopagebreak\vspace{-2.5pt}{\fontsize{8.9}{12.4}\selectfont\color{inksoft}\raggedright
  \noindent\llap{\leantag\hspace{9pt}}#1\par}\vspace{1pt}}
% a displayed statement in a tinted box: #1 colour, #2 text
\newcommand{\statement}[2]{\par\addvspace{7pt}\begingroup\setlength{\fboxsep}{8.5pt}\setlength{\fboxrule}{0.6pt}%
  \noindent\fcolorbox{#1!30!white}{#1!3!white}{\parbox{\dimexpr\linewidth-2\fboxsep-2\fboxrule\relax}{#2}}%
  \endgroup\par\addvspace{7pt}}
\newcommand{\nt}[2]{\par\noindent\hangindent=1.1em\hangafter=1 #1\enspace{\color{inksoft}#2}\par\vspace{1pt}}
\newcommand{\legdot}[1]{\tikz[baseline=-0.62ex]\fill[#1] (0,0) circle (2.5pt);\hspace{3pt}}
\newcommand{\crow}[4]{\csname CS@#1\endcsname & #2 & \textcolor{inksoft}{#3} & \textcolor{inksoft}{#4}\\}
\newcommand{\leanquoteline}[2]{\par\noindent\hangindent=\dimexpr#1\fontdimen2\font+2.2em\relax\hangafter=1
  \hspace*{#1\fontdimen2\font}#2}
\newcommand{\quoted}[1]{\par\vspace{3pt}{\ttfamily\fontsize{7.5}{10.4}\selectfont\color{leancol}\raggedright
  \setlength{\parskip}{0pt}\leanquote{#1}\par}\vspace{4pt}}
\newcommand{\bx}[3]{{\sffamily\bfseries\fontsize{7.3}{9.2}\selectfont\color{#1!82!black}#2}\\[1.5pt]#3}
\newcommand{\cs}[2]{\expandafter\newcommand\csname CS@#1\endcsname{#2}}

% theorems
\newtheoremstyle{plainx}{7pt plus 2pt}{7pt plus 2pt}{\itshape}{}{\bfseries}{.}{0.5em}{}
\newtheoremstyle{uprightx}{7pt plus 2pt}{7pt plus 2pt}{}{}{\bfseries}{.}{0.5em}{}
\theoremstyle{plainx}
\newtheorem{theorem}{Theorem}[section]
\newtheorem{lemma}[theorem]{Lemma}
\newtheorem{proposition}[theorem]{Proposition}
\newtheorem{corollary}[theorem]{Corollary}
\theoremstyle{uprightx}
\newtheorem{definition}[theorem]{Definition}
\newtheorem{setting}[theorem]{Setting}          % hypotheses shared by several statements; each of them cites it by number
\newtheorem{remark}[theorem]{Remark}
\newtheorem{example}[theorem]{Example}
\renewcommand{\qedsymbol}{{\setlength{\fboxsep}{0pt}\setlength{\fboxrule}{0.5pt}\fbox{\rule{0pt}{0.46em}\hspace{0.46em}}}}
% what holds throughout the note: the ambient objects only (implicit and instance arguments), never a hypothesis
\newenvironment{conventions}{\par\addvspace{5pt}\noindent\textbf{Conventions.}\enspace\ignorespaces}{\par\addvspace{3pt}}
\makeatletter
\renewenvironment{proof}[1][Proof]{\par\pushQED{\qed}\normalfont\topsep0pt\partopsep0pt
  \trivlist\item[\hskip\labelsep\itshape #1.]\ignorespaces}{\popQED\endtrivlist\@endpefalse}
\makeatother


% where each declaration of the concordance is treated: \cs{key}{place}, the place naming a \label or a lettered item
\cs{absnormbordereddet}{Lemma~\ref{l:bordered}}
\cs{existsliftslinkDiscSqle}{Proposition~\ref{p:lifts}}
\cs{existsisIntMatrix}{Lemma~\ref{l:int}}
\cs{existsorthogonalsplit}{Lemma~\ref{l:split}}
\cs{existsmemlinktermDisceq}{Lemma~\ref{l:vector}}
\cs{onelexRowsIndexandsaturationIndex}{Corollary~\ref{c:pos}}
\cs{minorNormsqlelinkDiscSqmul}{Proposition~\ref{p:minor}}
\cs{existslinksplitleclassGauge}{Theorem B}
\cs{pairInvariantmulle}{Theorem A}
% ----- macros-148: the recalled notions of unit 148 (chunks a, b) and of its neighbours, with every explicit
% ----- argument of the Lean definition as a slot. Shared by Notes F, P, T, G; copy unchanged into the preamble.
\newcommand{\fs}{\mathfrak{s}}                      % the shape (rho, kappa, k_0)
\newcommand{\sT}{\mathsf{T}}                        % a finite set of finite places
\newcommand{\cL}{\mathcal{L}}                       % a level
\newcommand{\AEf}{\mathbb{A}_{E,f}}
\newcommand{\hO}{\hat{O}_E}
\newcommand{\GL}{\mathrm{GL}}
\newcommand{\full}{\mathrm{full}}
\newcommand{\cut}{\mathrm{cut}}
\newcommand{\rh}[1]{\varrho(#1)}                    % Shape.ϱ
\newcommand{\kk}[1]{k(#1)}                          % Shape.k = #κ
\newcommand{\nn}[1]{n(#1)}                          % Shape.n = ϱ + 1
\newcommand{\rr}[1]{r(#1)}                          % Shape.r = 2ϱ + k + 1
\newcommand{\Ix}[1]{I(#1)}                          % Shape.BesselIndex
\newcommand{\Mx}[1]{M(#1)}                          % Shape.MiddleIndex
\newcommand{\ninf}[1]{f(#1)}                        % numInfPlaces E
\newcommand{\ET}[1]{E_{#1}}                         % LocProd T
% the components of the filling data S
\newcommand{\bS}[1]{b(#1)}                          % S.b
\newcommand{\NS}[1]{N(#1)}                          % S.N
\newcommand{\qS}[2]{q_{#2}(#1)}                     % S.levelPrime i
\newcommand{\gS}[1]{g(#1)}                          % S.g = q_0 q_1
\newcommand{\piS}[1]{\bar\pi(#1)}                   % S.piBar
\newcommand{\vpS}[1]{\varpi(#1)}                    % S.ϖ
\newcommand{\ChS}[2]{\mathrm{Ch}_{#2}(#1)}          % S.chosenPlaces i
\newcommand{\KcolS}[1]{K^{\mathrm{col}}(#1)}        % S.columnLevel
\newcommand{\BsysS}[2]{\mathfrak{B}_{#2}(#1)}       % S.besselSys v
% forms, groups, levels of the set-up (Setup.lean)
\newcommand{\Trat}[2]{T^{\mathrm{rat}}_{#1}(#2)}    % S.ratMiddle R
\newcommand{\Tlat}[2]{T^{\mathrm{lat}}_{#1}(#2)}    % S.latMiddle R
\newcommand{\GS}[2]{G_{#1}(#2)}                     % S.G R
\newcommand{\GFin}[1]{G^{\mathrm{lat}}_{#1}}        % S.GFin
\newcommand{\GRes}[1]{\mathsf{G}(#1)}               % S.GRes
\newcommand{\Cl}[1]{\mathrm{Cl}(#1)}                % S.Classes
\newcommand{\Bc}[1]{\mathrm{Bc}(#1)}                % S.BesselCoord
\newcommand{\KS}[1]{K(#1)}                          % S.latStabilizer
\newcommand{\redS}[1]{\mathrm{red}_{#1}}            % S.reduction
\newcommand{\KgS}[1]{K_g(#1)}                       % S.principalLevel
\newcommand{\redim}[2]{\overline{#2}^{\,#1}}        % S.reducedImage A
\newcommand{\kapS}[2]{\kappa_{#2}(#1)}              % S.kappa v
\newcommand{\gid}[1]{\hat{g}(#1)}                   % S.gIdele
\newcommand{\latB}[1]{\Lambda(#1)}                  % S.latBasis
\newcommand{\rinf}[1]{\mathrm{r}^{\infty}_{#1}}     % S.ratToInf
\newcommand{\rfin}[1]{\mathrm{r}^{f}_{#1}}          % S.ratToFin
\newcommand{\incH}[1]{\mathrm{inc}_{#1}}            % S.inclHFin
\newcommand{\hc}[2]{h_{#1}(#2)}                     % S.h c
\newcommand{\pos}[4]{\mathrm{pos}_{#1}(#2,#3;#4)}   % S.termPosition c c' γ
\newcommand{\terms}[4]{\mathrm{Tm}_{#1}(#2,#3,#4)}  % S.termsOver c s c'
\newcommand{\Zf}[2]{Z_{#1}(#2)}                     % S.frameProducts γ
\newcommand{\tdisc}[2]{\mathsf{g}_{#1}(#2)}         % S.termDisc γ
\newcommand{\clU}[2]{\mathcal{U}_{#1,#2}}           % S.classUnip c
\newcommand{\cutoff}[2]{\beta_{#1,#2}}              % S.classCutoff c
\newcommand{\besfn}[4]{W^{#1}_{#2}(#3)_{#4}}        % S.besselFn y i, with the side
\newcommand{\amp}[5]{a^{#1}_{#2}(#3,#4;#5)}         % S.amplitude c c' γ, with the side
\newcommand{\ampabs}[5]{|a|^{#1}_{#2}(#3,#4;#5)}    % S.amplitudeAbs c c' γ, with the side
\newcommand{\ampsum}[6]{A^{#1}_{#2}(#3;#4,#5,#6)}   % S.amplitudeSum T c s c', with the side
\newcommand{\uchar}[2]{\psi^{\mathcal{U}}_{#1}(#2)} % S.unipotentChar u
\newcommand{\ucov}[1]{\nu(#1)}                      % S.unipCovolume
\newcommand{\siegC}[1]{c_{*}(#1)}                   % S.siegelConst
\newcommand{\packC}[1]{a_{*}(#1)}                   % S.packetConst
\newcommand{\uord}[1]{\mathrm{uo}(#1)}              % S.unipOrder
\newcommand{\cdim}[1]{\mathrm{cd}(#1)}              % S.columnDim
\newcommand{\bdeg}[1]{\deg(#1)}                     % S.besselDegree
\newcommand{\mainE}[1]{\mathsf{d}(#1)}              % S.mainEnergy
\newcommand{\besE}[1]{\mathsf{d}_{1}(#1)}           % S.besselEnergy
\newcommand{\tailS}[1]{\mathrm{tl}(#1)}             % S.tailScale
\newcommand{\codimF}[1]{\mathrm{cf}(#1)}            % S.codimFactor
\newcommand{\fscale}[3]{\mathrm{fs}_{#1}(#2,#3)}    % S.frameScaleSq v a
\newcommand{\fsprod}[1]{\mathrm{fs}(#1)}            % S.frameScaleProd
\newcommand{\mass}[2]{\operatorname{mass}(#1,#2)}   % mass E r
% chunk b and the sieve
\newcommand{\energy}[3]{\mathrm{En}_{#1}(#2;#3)}    % Filling.energy A Ω (S implicit in Lean, shown)
\newcommand{\linkE}[5]{\mathrm{LE}^{#1}_{#2}(#3;#4,#5)} % Filling.linkEnergy T Ω' R (side and S implicit, shown)
\newcommand{\fixset}[3]{\mathrm{Fix}_{#1}(#2,#3)}   % Filling.fixingSet Sg d (S implicit, shown)
\newcommand{\vfix}[3]{\mathrm{Fix}_{#1,#2}(#3)}     % Filling.vectorFixer w d (S implicit, shown)
\newcommand{\charT}[2]{\psi_{#1,#2}}                % S.charOnPlaces T
\newcommand{\fgain}[3]{\beta_{#1,#2}(#3)}           % S.frequencyGain T ξ
\newcommand{\link}[1]{\operatorname{link}(#1)}      % link γ
\newcommand{\clink}[2]{\operatorname{link}^{\circ}(#1,#2)} % commonLink γ γ'
\newcommand{\tdnorm}[2]{\|\mathsf{g}\|_{#1}(#2)}    % S.termDiscNorm γ
\newcommand{\mnorm}[3]{\mathrm{mn}_{#1}(#2;#3)}     % S.minorNorm j γ
\newcommand{\cgauge}[5]{\mathrm{cg}_{#1}(#2;#3,#4;#5)} % S.classGauge j c c' γ
\newcommand{\pinv}[5]{\mathrm{pi}_{#1}(#2,#3;#4,#5)} % S.pairInvariant c c' γ γ'
\newcommand{\ldsq}[3]{\mathrm{ld}_{#1}(#2,#3)}      % S.linkDiscSq v e
\newcommand{\lmet}[1]{Q^{\mathrm{link}}_{#1}}       % S.linkMetric
% ----- end of macros-148
% ----- recalled notions of the set-up that this note needs beyond macros-148, each with every explicit argument
\newcommand{\Jrat}[1]{J_{#1}}                         % S.ratGram
\newcommand{\fpair}[1]{[\mathfrak{f}\mid #1\mathfrak{f}]}   % framePair γ (as in the notes on the set-up)
\newcommand{\ifp}[1]{[\mathfrak{e}\mid #1\mathfrak{e}]}     % isotropicFramePair γ
\newcommand{\lmat}[1]{\mathrm{L}(#1)}                 % linkMatrix γ
\newcommand{\zp}[1]{\mathrm{L}_{z}(#1)}               % zFormPair T
\newcommand{\lform}[1]{Q^{\pm}_{#1}}                  % S.linkForm
\newcommand{\cc}[2]{\mathrm{C}_{#1}(#2)}              % S.classCoord c
\newcommand{\ccX}[2]{\mathrm{C}^{\circ}_{#1}(#2)}     % S.classCoordX₀ c
\newcommand{\lifts}[4]{#4\in\operatorname{IL}(#2;#3)} % HasIntegralLifts c γ m (S is implicit in Lean; not shown, as in the set-up)
\newcommand{\xri}[3]{\mathrm{i}_{#1,#2}(#3)}          % S.xRowsIndex j k
\newcommand{\sat}[2]{\mathfrak{d}_{#1}(#2)}           % S.saturationIndex k
\newcommand{\ldisc}[3]{\operatorname{disc}_{#1,#2}(#3)} % S.latticeDisc c V
\newcommand{\fvi}[5]{\mathrm{o}_{#1,#2}(#3;#4,#5)}    % S.frameValueIndex c k a V
\newcommand{\haslevel}[3]{#3\in\operatorname{Lv}_{#2}(#1)} % S.HasLevel j γ
\newcommand{\lminor}[2]{D_{#2}(#1)}                   % leadingMinor Z j
\newcommand{\vol}{\operatorname{vol}}
\newcommand{\CC}{\mathbb{C}}
\newcommand{\RR}{\mathbb{R}}
\newcommand{\Agr}[2]{#2\in\operatorname{Agr}(#1)}     % S.Agrees s
\newcommand{\FixN}[2]{\operatorname{fix}(#1)\ge #2}   % S.FixesNormAtLeast s x

\begin{document}
\color{ink}

{\raggedright
{\sffamily\bfseries\fontsize{7.2}{9}\selectfont\color{leancol}\addfontfeatures{LetterSpace=8}READBACK NOTE 39\par}
\vspace{5pt}
{\fontsize{21.5}{25.5}\selectfont\bfseries Links of pairs of terms: integral lifts and the discriminant gain\par}
\vspace{5pt}
{\fontsize{12.4}{16}\selectfont\itshape\color{inksoft} One unit of the Lean artifact \ArtifactName, read back into
mathematics\par}}
\vspace{8pt}
{\color{rulegrey}\hrule height 0.6pt}
\vspace{7pt}
\thispagestyle{empty}

\textbf{Summary.} In the second moment of the Poincaré series, two rational terms $\gamma$, $\gamma'$ over the
same classes and the same element of the level group interact through their \emph{links}, two Lagrangian
subspaces of a split hermitian space of dimension \m{2(\kk{\fs}+1)}. This note proves the arithmetic input for
the pairs whose links differ: the product of the invariants of the pair, which measure the lattice of the
common part of the links, is smaller than the product of the gauges of the two terms by a power of \m{2/\gS{S}}
for each dimension in which the links differ. The proof computes the discriminant of the lattice of integral
lifts of a link by a bordered determinant and Jacobi's identity, splits this lattice orthogonally along the
second link, and applies an inequality of Liouville type to the pairing of the two complements, which is
divisible by \m{\gS{S}} because the two terms are congruent.

\statement{leancol}{\fontsize{9.6}{13.2}\selectfont\textbf{What this note is.} The artifact \ArtifactName{} is a
body of Lean code: the compact edition of a formal proof, due to L.~Alpöge with the AI system Claude, of the
Riemann hypothesis and of its generalisation to Dirichlet $L$-functions. That headline is a claim. The authors
report that every file compiles and that Lean's kernel accepted the proofs; an independent replay is recorded in
the handout's record of checks when it reports. No mathematician has read the argument through, and nothing here
says that the Riemann hypothesis is proved. This note (Note~P of its unit) treats \UnitInNote{} of the
\UnitMembers{} declarations of the unit \emph{Translates of Poincaré series fill the rows}: the module on link
pairings. Proved here from the inputs: Theorems A and~B and the statements of §\ref{s:proofs}; where a step
is a verification of matrix entries, the note says so and gives the identity that is verified. Taken for
granted and stated as Inputs 1 to~4: the properties of links, levels and class coordinates from the set-up of
the unit, an inequality of Liouville type for Gram determinants, the bound of an index by agreeing characters,
and Jacobi's identity for discriminants of complements from the background library. Every statement was
compared with its Lean text, argument by argument. The Lean text read is \PinText; Appendix~A is the concordance
of \ConcDecls{} declarations.}

\sect{The statements}{s:intro}
The objects are recalled in §\ref{s:setting}. In brief: \m{\link{\gamma}} is the link of a rational element,
\m{\clink{\gamma}{\gamma'}} the common part of two links inside~$E^\kappa$; \m{\pinv{S}{c}{c'}{\gamma}{\gamma'}} is the
invariant of the pair and \m{\cgauge{S}{j}{c}{c'}{\gamma}} the gauge of a term of level~$j$.

\statement{arith}{\phantomsection\label{t:A}\textbf{Theorem A} (the discriminant gain).\enspace
\emph{In Setting~\ref{s:S}, suppose that $E$ is stably unramified. Let $c$, $c'$ be classes of~$S$,
\m{s\in\GRes{S}} an element that agrees, $x$ a natural number such that $s$ fixes norm at least~$x$, $j$ a
level, and \m{\gamma,\gamma'\in\terms{S}{c}{s}{c'}} two terms with \m{\haslevel{S}{j}{\gamma}} and
\m{\haslevel{S}{j}{\gamma'}}. Then there is a natural number $i$ such that}
\begin{enumerate}[label=\textup{(\roman*)},leftmargin=2.3em,itemsep=1pt]
\item \m{\kk{\fs}\le\dim_E\clink{\gamma}{\gamma'}+i+(j-\rh{\fs})_+};
\item \emph{if \m{i=0}, then \m{\link{\gamma}=\link{\gamma'}}; and}
\item \emph{with \m{x'=x} if \m{\rh{\fs}<j} and \m{i=0}, and \m{x'=1} otherwise,}
\[\begin{aligned}&\pinv{S}{c}{c'}{\gamma}{\gamma'}\cdot\pinv{S}{c}{c'}{\gamma'}{\gamma}\cdot x'^2\\
&\qquad\le\ \Bigl(\frac{2}{\gS{S}}\Bigr)^{2\,i\,\ninf{E}}\cdot\cgauge{S}{j}{c}{c'}{\gamma}\cdot\cgauge{S}{j}{c}{c'}{\gamma'} .\end{aligned}\]
\end{enumerate}}
\leanline{\lean{pairInvariantmulle}}

\statement{arith}{\phantomsection\label{t:B}\textbf{Theorem B} (splitting a link along a second one).\enspace
\emph{In Setting~\ref{s:S}, suppose that $E$ is stably unramified. Let $c$, $c'$ be classes of~$S$,
\m{s\in\GRes{S}} an element that agrees, $x$ a natural number such that $s$ fixes norm at least~$x$, $j$ a
level, \m{\gamma\in\terms{S}{c}{s}{c'}} a term with \m{\haslevel{S}{j}{\gamma}}, and \m{\gamma'\in\GS{S}{E}}. Put
\m{k_\gamma=\pos{S}{c}{c'}{\gamma}} and \m{L=\link{\gamma}\cap\link{\gamma'}}. Then}
\begin{enumerate}[label=\textup{(\roman*)},leftmargin=2.3em,itemsep=1pt]
\item \m{0<\ldisc{S}{c}{\clink{\gamma}{\gamma'}}} \emph{and} \m{1\le\fvi{S}{c}{k_\gamma}{1}{\clink{\gamma}{\gamma'}}};
\item \emph{there are a natural number~$i$, a finite index set $\mathcal C$, a matrix $m$ over $E$ with
columns indexed by~$\mathcal C$, a matrix $y$ over $E$ with $i$ columns, \m{v\in\GL_i(\AEf)} and a matrix $z$
over \m{\AEf} with rows $\mathcal C$ and $i$ columns, such that: \m{i=\kk{\fs}+1-\dim L};
\m{\kk{\fs}\le\dim\clink{\gamma}{\gamma'}+i+(j-\rh{\fs})_+}; if \m{i=0} then \m{L=\link{\gamma}}; the columns of
$m$ span~$L$; the columns of $y$ lie in \m{\link{\gamma}}; \m{y\,t\in L} only for \m{t=0}; the columns of
\m{(m\,|\,y)} span \m{\link{\gamma}}; \m{\lifts{S}{c}{\gamma}{\tilde mz+\tilde yv}}; and, with $x'$ as in
Theorem~A,}
\[\bigl(\pinv{S}{c}{c'}{\gamma}{\gamma'}\cdot x'\bigr)^2\cdot\ldsq{S}{y}{v}\ \le\ \cgauge{S}{j}{c}{c'}{\gamma}^2 .\]
\end{enumerate}}
\leanline{\lean{existslinksplitleclassGauge}}

Theorem~A is used in the crude estimate of the unit (the note on the energy of amplitude sums): there the pairs
of terms of one level over one \m{(c,s,c')} are summed with the weights that the invariants of the pair
control, and the factor \m{(2/\gS{S})^{2i\ninf{E}}} pays for the number of links at distance~$i$.

\sect{The setting}{s:setting}
\begin{conventions}\phantomsection\label{c:conventions}
The natural numbers are \m{0,1,2,\dots}, and \m{(a)_+} is $a$ if \m{a\ge0} and $0$ otherwise. $E$ is a CM field
with its CM type: a number field that is a CM field, with complex conjugation \m{x\mapsto\bar x} and one
embedding at each infinite place~$\tau$; \m{x_\tau} is the image of~$x$, \m{\ninf{E}} the number of infinite
places, \m{N_{E/\Q}} the norm. \m{\AEf} is the ring of finite adeles, \m{\hO} its subring of integral adeles,
\m{\tilde a} the image in \m{\AEf} of a matrix $a$ over~$E$, and \m{[\xi]} the generalised index of
\m{\xi\in\AEf}: the index of \m{\xi\hO} in \m{\hO} for integral~$\xi$, extended multiplicatively to the units
and $0$ on non-units; for \m{\xi\in E^\times} it is \m{|N_{E/\Q}(\xi)|}. $E$ is \emph{stably unramified} if
it contains a primitive root of unity of order~$5040$. A \emph{shape} is a triple
\m{\fs=(\varrho,\kappa,k_0)}; \m{\rh{\fs}}, \m{\kk{\fs}=\#\kappa}; a \emph{level} is an element $j$ of
\m{\{0,\dots,\rh{\fs}+1\}}. For a matrix $X$ over a ring with involution, $X^*$ is the conjugate transpose; a
matrix \emph{has a left inverse} if its columns are linearly independent, and its \emph{column span} is the
space spanned by its columns. Index sets written $\mathcal C$, $\mathcal C'$ are finite index sets of columns
of matrices; where the formal text lists such a set, its finiteness or a matrix among the arguments of a
statement that does not use them, this is said in §\ref{s:notes}. As in Note~F, and unlike in the handout's
overview, $c$, $c'$ are classes of the filling data and $s$, $t$ elements of the finite group \m{\GRes{S}}; the
complex variable and the unit class of the overview do not occur.
\end{conventions}

\begin{setting}[filling data]\label{s:S}
$E$ is a CM field with its CM type, $\fs$ is a shape, and $S$ is filling data over $E$ of shape~$\fs$, in the
sense of the set-up of this unit, with its prime \m{\NS{S}}, the number \m{\gS{S}}, the finite group
\m{\GRes{S}}, the classes \m{\Cl{S}}, the groups \m{\GS{S}{E}} and \m{\GFin{S}}, the integral subgroup \m{\KS{S}}
with its reduction \m{\redS{S}} and kernel \m{\KgS{S}}, the positions \m{\pos{S}{c}{c'}{\gamma}} and the sets of
terms \m{\terms{S}{c}{s}{c'}}, as recalled in the setting ``filling data'' of Note~F of this unit (\emph{Translates
of Poincaré series fill the rows}).
\end{setting}

\subhead{Recalled definitions.} All the notions below belong to the set-up of the unit; they are recalled with
every argument of their formal definitions. The index set \m{\Ix{\fs}} of the whole space consists of
\m{e_0,\dots,e_{\rh{\fs}-1}}, of~$\kappa$, of $z$ and of \m{e^*_0,\dots,e^*_{\rh{\fs}-1}}; \m{\Mx{\fs}=\kappa\sqcup\{z\}}
indexes the \emph{middle}. $X$ is the span of the~$e_a$; the frame is \m{\mathfrak f_a=e_a} (\m{a<\rh{\fs}}),
\m{\mathfrak f_{\rh{\fs}}=z}. In Setting~\ref{s:S}: \m{T=\Trat{S}{E}} is the diagonal matrix on \m{\Mx{\fs}} with
entries $1$ on $\kappa$ and \m{-\NS{S}} at~$z$; \m{\Jrat{S}} is its split Gram matrix on \m{\Ix{\fs}}, the Gram
matrix of the hermitian form $h$; \m{\Tlat{S}{\AEf}} and its split Gram matrix $J_f$ are the same in the lattice
basis over the finite adeles.
\begin{enumerate}[label=\textup{(\alph*)},leftmargin=2.3em,itemsep=2pt]
\item For \m{\gamma\in\GS{S}{E}}: $\Zf{S}{\gamma}=(h(\mathfrak f_a,\gamma\mathfrak f_{a'}))_{a,a'}$ is the frame Gram
matrix, \m{\lminor{Z}{l}} the leading minor of order $l$ of a matrix~$Z$, and the \emph{term discriminant} is
\[\tdisc{S}{\gamma}=|\lminor{\Zf{S}{\gamma}}{\varrho+1}|^2-\NS{S}^2|\lminor{\Zf{S}{\gamma}}{\varrho}|^2,\]
with \m{|\xi|^2=\xi\bar\xi}; \m{\tdnorm{S}{\gamma}=\prod_\tau|\tdisc{S}{\gamma}_\tau|}.
\item \m{\fpair{\gamma}=(\mathfrak f\,|\,\gamma\mathfrak f)} is the matrix of the frame and its image,
\m{\ifp{\gamma}=(e_a\,|\,\gamma e_a)_{a<\rh{\fs}}} its part on~$X$. The \emph{link matrix} \m{\lmat{\gamma}} sends a
vector $t$ to the pair of the middle parts of $t$ and of \m{\gamma^{-1}t}, a vector indexed by
\m{\Mx{\fs}\sqcup\Mx{\fs}}; and \m{\link{\gamma}=\lmat{\gamma}\bigl(\ker(\ifp{\gamma}^*\Jrat{S})\bigr)}, the
\emph{link}: the pairs of the classes of $t$ and of $\gamma^{-1}t$ for $t$ orthogonal to $X$ and~$\gamma X$.
\item \m{\zp{T}} is the matrix with two rows that takes from a pair of middle vectors the two products with
the row of $T$ at~$z$; \m{\lform{S}=\diag(T,-T)} is the \emph{link form} and
\[\lmet{S}=\tfrac12\diag(T,T)+\NS{S}^{-1}\,\zp{T}^*\zp{T}\]
the \emph{link metric}, both on pairs of middle vectors. For \m{\gamma,\gamma'}, the subspace
\m{\clink{\gamma}{\gamma'}} of $E^\kappa$ is the space of the $\kappa$-parts of the first components of the
pairs in $\link{\gamma}\cap\link{\gamma'}\cap\ker\zp{T}$.
\item For a matrix $v$ over $E$ on pairs of middle vectors with columns~$\mathcal C$ and
\m{e\in\GL_{\mathcal C}(\AEf)}: \m{\ldsq{S}{v}{e}=|N_{E/\Q}\det(v^*\lmet{S}v)|\cdot[\det e]^2}.
\item For a class~$c$, \m{\cc{S}{c}} is the matrix over \m{\AEf} that takes rational coordinates on the whole
space to coordinates in the basis of the lattice of~$c$, and \m{\ccX{S}{c}} the matrix whose columns are the
basis vectors of $E^\kappa$ in the basis of the middle part of that lattice. For classes $c$, $c'$,
\m{\gamma\in\GS{S}{E}} and \m{k=\pos{S}{c}{c'}{\gamma}}: \m{\cc{S}{c}\,\widetilde{\fpair{\gamma}}=\fpair{k}}.
\item For a class $c$, \m{\gamma\in\GS{S}{E}} and a matrix $m$ over \m{\AEf} on pairs of middle vectors with
columns~$\mathcal C$: \m{\lifts{S}{c}{\gamma}{m}} (\emph{$m$ has integral lifts}) if there are a finite index set
$\mathcal C'$, a matrix $t$ over $E$ with rows \m{\Ix{\fs}} and columns $\mathcal C'$ and a matrix $w$ over
\m{\AEf} with rows $\mathcal C'$ and columns $\mathcal C$ such that \m{\ifp{\gamma}^*\Jrat{S}\,t=0}, the matrix
\m{\cc{S}{c}\,\tilde t\,w} is integral, and \m{\widetilde{\lmat{\gamma}t}\,w=m}.
\item For \m{k\in\GFin{S}} with frame Gram matrix $Z(k)$ for \m{\Tlat{S}{\AEf}} and a level~$j$:
\m{\xri{S}{j}{k}} is the index of the \m{\hO}-span of the conjugated first \m{\min(j,\rh{\fs})} rows of $Z(k)$ in
that of its conjugated first \m{\rh{\fs}} rows, and \m{\sat{S}{k}} the index of the \m{\hO}-span of the columns of
\m{\ifp{k}} in the lattice of the integral vectors of their span; both are $0$ if infinite.
\item For a class $c$ and a subspace $V$ of~$E^\kappa$: \m{\ldisc{S}{c}{V}} is the discriminant of the lattice of
the adelic points of $V$ that are integral in the lattice of~$c$; for \m{k\in\GFin{S}} and a unit $a$ of
\m{\AEf}, \m{\fvi{S}{c}{k}{a}{V}} is the index in the lattice of the values \m{(h(k\mathfrak f_l,y))_{l<\rh{\fs}}}, for $y$
in that lattice of~$V$, of its intersection with $a$ times the lattice spanned by the conjugated first
\m{\rh{\fs}} rows of $Z(k)$ ($0$ if infinite). The \emph{invariant of a pair} is
\[\begin{aligned}\pinv{S}{c}{c'}{\gamma}{\gamma'}={}&\ldisc{S}{c}{\clink{\gamma}{\gamma'}}\cdot
\fvi{S}{c}{\pos{S}{c}{c'}{\gamma}}{1}{\clink{\gamma}{\gamma'}}\\&\cdot\sat{S}{\pos{S}{c}{c'}{\gamma}} .\end{aligned}\]
\item For a level $j$, \m{\gamma\in\GS{S}{E}} and classes $c$, $c'$, with \m{k=\pos{S}{c}{c'}{\gamma}}, the
\emph{minor norm} and the \emph{gauge} are
\[\mnorm{S}{j}{\gamma}=|N_{E/\Q}(\lminor{\Zf{S}{\gamma}}{j})|\cdot\NS{S}^{-2(j-\rh{\fs})_+\ninf{E}},\]
\[\cgauge{S}{j}{c}{c'}{\gamma}=\frac{\max\bigl(\tdnorm{S}{\gamma},\ \NS{S}^{\ninf{E}}\mnorm{S}{j}{\gamma}\bigr)}
{\xri{S}{j}{k}\cdot\xri{S}{j}{k^{-1}}} .\]
\item \m{\haslevel{S}{j}{\gamma}} (\emph{$\gamma$ has level $j$}) if $\gamma$ is in normal form of index~$j$,
\m{\lminor{\Zf{S}{\gamma}}{j}\ne0}, and \m{\rh{\fs}<j} holds exactly if \m{\tdisc{S}{\gamma}\ne0}. Here $\gamma$ is in
\emph{normal form of index $j$} if there is \m{\upsilon\in E} with \m{\upsilon\bar\upsilon=1} such that for every frame
index \m{a\ge j} there are coefficient vectors $x$, $y$ supported on the indices \m{\le a} with \m{x_a=\upsilon},
\m{y_a=1} and \m{\sum_{a'}x_{a'}\mathfrak f_{a'}=\gamma\sum_{a'}y_{a'}\mathfrak f_{a'}}.
\item \m{s\in\GRes{S}} \emph{agrees}, in symbols \m{\Agr{S}{s}}, if for all \m{u,u'\in\KS{S}} in the Bessel
unipotent group with $\redS{S}(u)\,s=s\,\redS{S}(u')$ the unipotent characters of $u$ and $u'$ are equal. For a
natural number~$x$, $s$ \emph{fixes norm at least}~$x$, in symbols \m{\FixN{s}{x}}, in the sense of the
definition of fixing sets recalled in Note~F: there is a finite set of triples of an index \m{i\in\{0,1\}}, a
chosen place above the $i$-th prime of $S$ and a Boolean, the product of whose primes is at least~$x$, such
that for each of them $s$, or its inverse adjoint if the Boolean is true, fixes at that place a non-zero vector
of the span of the frame.
\end{enumerate}
The symbols of (a) to~(k) are those of the notes on the set-up of this unit, which state the definitions; the
predicates of~(k) are written there with their symbols and here in words.

\sect{What is taken for granted}{s:inputs}
\statement{est}{\phantomsection\label{i:1}\textbf{Input 1} (links, levels, class coordinates; set-up of this
unit).\enspace In Setting~\ref{s:S}, for \m{\gamma,\gamma'\in\GS{S}{E}}:
\begin{enumerate}[label=\textup{(\alph*)},leftmargin=2.3em,itemsep=1pt]
\item \m{\dim_E\link{\gamma}=\kk{\fs}+1}, and \m{\link{\gamma}} is Lagrangian for the link form: a pair $p$ lies
in \m{\link{\gamma}} exactly if \m{q^*\lform{S}\,p=0} for all \m{q\in\link{\gamma}}. In particular
\m{v^*\lform{S}u=0} for matrices $u$, $v$ with columns in \m{\link{\gamma}}.
\item If the columns of $t$ are orthogonal to \m{X+\gamma X} and those of $t'$ to \m{X+\gamma'X}, then for
\m{a,a'\in E}:
\m{(\lmat{\gamma}t)^*\diag(aT,-a'T)(\lmat{\gamma'}t')=a\,t^*\Jrat{S}t'-a'\,t^*\Jrat{S}(\gamma\gamma'^{-1})t'}.
\item The link metric is positive definite at every infinite place, and \m{\det(v^*\lmet{S}v)\ne0} for every
matrix $v$ with a left inverse.
\item For matrices $t$, $t'$ over $E$ with rows \m{\Ix{\fs}}:
\m{(\cc{S}{c}\tilde t)^*J_f(\cc{S}{c}\tilde t')=\widetilde{t^*\Jrat{S}t'}}; $J_f$ is unimodular; and the frame
Gram matrix of \m{\pos{S}{c}{c'}{\gamma}} is the image of~\m{\Zf{S}{\gamma}}.
\item If \m{\haslevel{S}{j}{\gamma}}, then \m{\haslevel{S}{j}{\gamma^{-1}}} and \m{\mnorm{S}{j}{\gamma}>0}. If
$\gamma$ is in normal form of index~$j$, then for every vector $y$ of frame coordinates there are $x$, $y'$
with \m{y'_a=y_a} for \m{a\ge j} and \m{\sum_ax_a\mathfrak f_a=\gamma\sum_ay'_a\mathfrak f_a}; in such a relation the
top coordinates satisfy \m{\bar x_\varrho x_\varrho=\bar y'_\varrho y'_\varrho}. If \m{\haslevel{S}{j}{\gamma}} and \m{b\ge j} is
a frame index, the combination of the vectors \m{\gamma\mathfrak f_a} with the coefficients of column $b$ of the
inverse of the matrix \m{\mathrm{pl}_j(\Zf{S}{\gamma})}, whose rows \m{a<j} are those of \m{\Zf{S}{\gamma}} and whose
other rows are those of the identity, lies in the span of the frame.
\item If \m{\Zf{S}{\gamma}\,y=D\,e_\varrho} for a coefficient vector $y$ and \m{D\in E}, and
\m{t=\sum_ay_a\gamma\mathfrak f_a}, then the coordinates of $t$ at the $e^*_a$ vanish, \m{-\NS{S}\,t_z=D}, and
\m{\NS{S}\sum_{i\in\kappa}|t_i|^2=|D|^2-\NS{S}^2|y_\varrho|^2}.
\item If \m{\gamma\in\terms{S}{c}{s}{c'}} and \m{\gamma'\in\terms{S}{c'}{s'}{c''}}, then
\m{\gamma\gamma'\in\terms{S}{c}{ss'}{c''}}; the inverse of a term over \m{(c,s,c')} is a term over
\m{(c',s^{-1},c)}; and the position of a term over \m{(c,1,c)} lies in \m{\KgS{S}}, so that it is congruent to
$1$ modulo \m{\gS{S}}.
\item If the columns of $n$ are orthogonal to \m{X+kX} for \m{k\in G(\Tlat{S}{\AEf})} and
\m{\zp{\Tlat{S}{\AEf}}\lmat{k}n=0}, then \m{\mathrm{fm}(k)^*\,\Tlat{S}{\AEf}\,n_M=-Z_X(k)^*\,n_X}, where $n_X$ is the
part of $n$ on the~$e_a$, $n_M$ its middle part, \m{\mathrm{fm}(k)} the matrix of the middle parts of the
\m{k\mathfrak f_l}, \m{l<\rh{\fs}}, and \m{Z_X(k)} the matrix of the first \m{\rh{\fs}} rows of $Z(k)$.
\end{enumerate}}

\statement{est}{\phantomsection\label{i:2}\textbf{Input 2} (Liouville's inequality for Gram determinants;
set-up of this unit).\enspace Let $E$ be stably unramified, $Q$ a square matrix over $E$ that is positive
definite at every infinite place, $a$, $a'$ matrices over $E$ with the rows of $Q$ and columns indexed by a
finite set $\mathcal C$ with \m{\det(a^*Qa')\ne0}, \m{v,v'\in\GL_{\mathcal C}(\AEf)}, and \m{x_0\ge0} a real
number such that \m{x_0\le[\delta]} whenever \m{\delta=\det(v^*\,\widetilde{a^*Qa'}\,v')} is a unit of~\m{\AEf}.
Then \m{x_0^2\le[\det v]^2\,|N_{E/\Q}\det(a^*Qa)|\cdot|N_{E/\Q}\det(a'^*Qa')|\,[\det v']^2}.}

\statement{est}{\phantomsection\label{i:3}\textbf{Input 3} (lattices; set-up of this unit and background
library).\enspace Let $E$ be stably unramified.
\begin{enumerate}[label=\textup{(\alph*)},leftmargin=2.3em,itemsep=1pt]
\item \emph{(Primitive completion.)} If a matrix over \m{\AEf} with columns \m{\mathcal C\sqcup\mathcal C'} has a
left inverse, it becomes primitive (its columns extend to a basis of the lattice of integral vectors) after
multiplication on the right by a block matrix with diagonal blocks \m{a\in\GL_{\mathcal C}(\AEf)},
\m{e\in\GL_{\mathcal C'}(\AEf)}, an upper right block $z$ and lower left block~$0$. A square primitive matrix is
integral with integral inverse, so its determinant has generalised index~$1$.
\item \emph{(Jacobi's identity for discriminants of complements; reference to be confirmed.)} Let $J$ be a
unimodular hermitian matrix over~\m{\AEf}, $F$ an integral matrix, $n$ a primitive matrix with \m{F^*Jn=0}, of
complementary sizes, $p$ a matrix of bordering columns and $c_0$ a unit. Then the solutions $u'$ of
\m{n^*Ju'=0} are the combinations of the columns of~$F$, and
\[\bigl[\det\bigl(n^*(J+c_0^{-1}Jpp^*J)n\bigr)\bigr]\cdot\mathfrak i(F)^2\cdot[c_0]^{d_p}=[\det M],\]
where $M$ is the Gram matrix of \m{(F\,|\,p)} for $J$ with $c_0$ added on the diagonal block of~$p$, $d_p$ is
the number of columns of~$p$, and \m{\mathfrak i(F)} the index of the lattice of integral vectors in the lattice
of the $y$ with $Fy$ integral.
\item \emph{(Extension of a basis modulo a kernel.)} Let $A$ and $x$ be matrices over $E$ with
\m{\det(Ax)\ne0}, and $L$ a subspace containing the columns of~$x$. Then there is a matrix $y$ with
\m{\dim(L\cap\ker A)} columns such that \m{(x\,|\,y)} has a left inverse and its columns span~$L$, the columns
of $y$ lie in~$L$, \m{Ay=0}, and \m{yt} lies in the column span of $x$ only for \m{t=0}.
\item \emph{(Discriminant of a sublattice.)} Let $c$ be a class, $V$ a subspace of~$E^\kappa$ of
dimension~$d$, $p$ a matrix over $E$ with $d$ columns in~$V$, \m{a\in\GL_d(\AEf)} with
\m{D=[\det(p^*p)]\,[\det a]^2\ne0}, and $k$ in \m{\GFin{S}}. If \m{\ccX{S}{c}\,\tilde p\,a} is integral and its
values \m{(h(k\mathfrak f_l,\cdot))_l} lie in the lattice spanned by the conjugated first \m{\rh{\fs}} rows of
$Z(k)$, then \m{(\ldisc{S}{c}{V}\cdot\fvi{S}{c}{k}{1}{V})^2\le D}, \m{\ldisc{S}{c}{V}>0} and
\m{\fvi{S}{c}{k}{1}{V}\ge1}.
\end{enumerate}}

\statement{est}{\phantomsection\label{i:4}\textbf{Input 4} (agreeing characters force a large index; set-up
of this unit).\enspace In Setting~\ref{s:S}, let $E$ be stably unramified, \m{k\in\KS{S}} with \m{\redS{S}(k)}
agreeing, $x$ a natural number such that \m{\redS{S}(k)} fixes norm at least~$x$, and $n$ an integral matrix
with rows \m{\Ix{\fs}} and two columns whose \emph{pairings} with the $e_a$ and the $ke_a$ vanish, that is
\m{\ifp{k}^*J_f\,n=0}; here the pairing of a column $v$ of $n$ with a vector $w$ is \m{w^*J_fv=h(w,v)}. If the
determinant $\delta$ of the \m{2\times2} matrix of the pairings of the columns of $n$ with $z$ and~$kz$ is a
unit of~\m{\AEf}, then \m{x\le[\delta]}. (The formal statement is about an integral matrix $t$ and the plain
products \m{\fpair{k}^*t}; it is used for \m{t=J_fn}, which is integral because $n$ and $J_f$ are.)}

\sect{Proofs}{s:proofs}
Throughout, in Setting~\ref{s:S}, \m{N=\NS{S}}, \m{\varrho=\rh{\fs}}, \m{T=\Trat{S}{E}}, and \m{[N]=N^{2\ninf{E}}}
is the generalised index of the integer~$N$, the degree of $E$ being \m{2\ninf{E}}.

\subhead{The bordered determinant.}
\begin{lemma}[a bordered Gram determinant]\label{l:bordered}
In Setting~\ref{s:S}, let $j$ be a level and \m{\gamma\in\GS{S}{E}} with \m{\haslevel{S}{j}{\gamma}}. Let $P$ be
the matrix over $E$ with rows \m{\Ix{\fs}} whose columns are: the vectors $e_a$ for \m{a<\rh{\fs}}; the vectors
\m{\gamma e_a} for \m{a<\rh{\fs}} with \m{a<j}; the basis vectors $e^*_a$ for \m{a<\rh{\fs}} with \m{a\ge j}; and
$z$ and~$\gamma z$. Let $M$ be the sum of \m{P^*\Jrat{S}P} and of the diagonal matrix with entry \m{\NS{S}} at
the columns $e^*_a$, $z$, $\gamma z$ and $0$ at the others. Then
\[[\det\tilde M]=\bigl(\NS{S}^{2\ninf{E}}\cdot\mnorm{S}{j}{\gamma}\bigr)^2 .\]
\end{lemma}
\begin{proof}
Let \m{Z'=\Zf{S}{\gamma^{-1}}}, the conjugate transpose of \m{\Zf{S}{\gamma}}, and let $A$ be the square matrix of
size \m{\varrho+1} whose rows \m{a<j}, \m{a<\varrho}, are those of~$Z'$, whose rows \m{j\le a<\varrho} are those
of the identity, and whose top row is that of~$Z'$.

Order the columns of $P$ as the frame \m{\mathfrak f_0,\dots,\mathfrak f_\varrho} (the $e_a$, \m{a<\varrho}, and~$z$),
followed by the \m{\varrho+1} columns \m{\gamma e_b} (\m{b<\min(j,\varrho)}), $e^*_b$ (\m{j\le b<\varrho}) and
\m{\gamma z}, which we index by \m{b\le\varrho}, the column \m{\gamma z} having the index~$\varrho$. On the frame,
$M$ is the matrix \m{(h(\mathfrak f_a,\mathfrak f_{a'}))_{a,a'}+\diag(0,\dots,0,N)}, which is zero: the $e_a$ are
isotropic and orthogonal to~$z$, and \m{h(z,z)=-N}. The block of the frame against the other columns has the
entries
\[h(\mathfrak f_a,\gamma e_b)=\Zf{S}{\gamma}_{ab},\qquad h(\mathfrak f_a,e^*_b)=\delta_{ab},\qquad
h(\mathfrak f_a,\gamma z)=\Zf{S}{\gamma}_{a\varrho};\]
its columns \m{b<\min(j,\varrho)} and \m{b=\varrho} are those of \m{\Zf{S}{\gamma}=Z'^*} and its columns
\m{j\le b<\varrho} are those of the identity, so it is~$A^*$. As $M$ is hermitian,
\[M=\begin{pmatrix}0&A^*\\A&*\end{pmatrix},\qquad\det M=\pm\det A\cdot\overline{\det A},\]
and as \m{\det M\in E}, $[\det\tilde M]=|N_{E/\Q}(\det A)|^2$. It remains to show
$|N_{E/\Q}(\det A)|=N^{2\ninf{E}}\mnorm{S}{j}{\gamma}$.

If \m{j=\varrho+1}, then \m{A=Z'}, \m{|N_{E/\Q}\det Z'|=|N_{E/\Q}\lminor{\Zf{S}{\gamma}}{\varrho+1}|}, and the claim
is the definition of \m{\mnorm{S}{\varrho+1}{\gamma}}, in which \m{(j-\varrho)_+=1}.

Let \m{j\le\varrho}. The element \m{\gamma^{-1}} has level~$j$ (Input~1(e)), so for the frame index
\m{\varrho\ge j} the combination of the \m{\gamma^{-1}\mathfrak f_a} with the coefficients $y$ of the top column of
\m{\mathrm{pl}_j(Z')^{-1}} lies in the span of the frame: $\sum_ay_a\gamma^{-1}\mathfrak f_a=\sum_ax_a\mathfrak f_a$.
The matrix $A$ differs from \m{\mathrm{pl}_j(Z')} only in the top row, so by Cramer's rule
\m{\det A=\det\mathrm{pl}_j(Z')\cdot(Z'y)_\varrho}, and \m{\det\mathrm{pl}_j(Z')=\lminor{Z'}{j}\ne0}. Now
$(Z'y)_\varrho=h(z,\sum_ay_a\gamma^{-1}\mathfrak f_a)=h(z,\sum_ax_a\mathfrak f_a)=-N\,x_\varrho$, and
\m{\bar x_\varrho x_\varrho=\bar y_\varrho y_\varrho=1}, because the top diagonal entry of
\m{\mathrm{pl}_j(Z')\,\mathrm{pl}_j(Z')^{-1}} is~$1$ and the top row of \m{\mathrm{pl}_j(Z')} is that of the identity.
So \m{\det A=-N\,x_\varrho\,\lminor{Z'}{j}} with $x_\varrho$ of norm one at every place, and
\[|N_{E/\Q}(\det A)|=N^{2\ninf{E}}\,|N_{E/\Q}\lminor{\Zf{S}{\gamma}}{j}|=N^{2\ninf{E}}\mnorm{S}{j}{\gamma},\]
the leading minors of $Z'$ being the conjugates of those of \m{\Zf{S}{\gamma}}.
\end{proof}
\leanline{\lean{absnormbordereddet}}

\subhead{Integral lifts of a basis of the link.}
\begin{proposition}[lifts with small discriminant]\label{p:lifts}
In Setting~\ref{s:S}, suppose that $E$ is stably unramified. Let $c$, $c'$ be classes of~$S$, \m{t\in\GRes{S}},
$j$ a level, \m{\gamma\in\terms{S}{c}{t}{c'}} with \m{\haslevel{S}{j}{\gamma}}, $\mathcal C$ a finite index set and
$v$ a matrix over $E$ with rows \m{\Mx{\fs}\sqcup\Mx{\fs}} and columns $\mathcal C$ that has a left inverse and
whose column span is \m{\link{\gamma}}. Put \m{k=\pos{S}{c}{c'}{\gamma}} and
\m{w_0=\xri{S}{j}{k}\cdot\xri{S}{j}{k^{-1}}\cdot\sat{S}{k}}. Then there is \m{e\in\GL_{\mathcal C}(\AEf)} such that
\[\lifts{S}{c}{\gamma}{\tilde v\,e},\qquad\ldsq{S}{v}{e}\cdot w_0^2\le\mnorm{S}{j}{\gamma}^2,\qquad w_0\ne0 .\]
\end{proposition}
\begin{proof}
Let $B$ be the set of the \m{a<\varrho} with \m{a\ge j}; it is empty if \m{j\ge\varrho}.

\emph{Step 1: a rational lift.} There is a matrix $u$ over $E$ with rows \m{\Ix{\fs}} and columns
\m{B\sqcup\mathcal C} such that the columns of $u$ are orthogonal to \m{X+\gamma X}; \m{\lmat{\gamma}u=(0\,|\,v)};
and the $e_a$-coordinates of the columns of $u$, for \m{a\in B}, form the matrix \m{(1\,|\,0)}. Indeed, a
column $v_b$ of $v$ lies in the link, so \m{v_b=\lmat{\gamma}t_b} with $t_b$ orthogonal to \m{X+\gamma X}. By the
normal form of \m{\gamma^{-1}} (Input~1(e)), for every prescription of coordinates at the indices \m{a\in B}
there is a vector \m{w'=\sum_ay'_ae_a} with these coordinates, with \m{y'_\varrho=0}, such that
\m{\gamma^{-1}w'=\sum_ax_a\mathfrak f_a} with \m{x_\varrho=0}. Then $w'$ lies in \m{X\cap\gamma X}; it is orthogonal
to \m{X+\gamma X}, both spaces being isotropic, and \m{\lmat{\gamma}w'=0}, the middle parts of $w'$ and of
\m{\gamma^{-1}w'} being zero. Adding such vectors to the $t_b$, and taking them alone for the columns of~$B$,
gives~$u$.

\emph{Step 2: completion.} The matrix \m{\cc{S}{c}\tilde u} has a left inverse. By Input~3(a) there are
\m{a\in\GL_B(\AEf)}, \m{e\in\GL_{\mathcal C}(\AEf)} and $z_1$ such that
\m{n=\cc{S}{c}\,\tilde u\,\bigl(\begin{smallmatrix}a&z_1\\0&e\end{smallmatrix}\bigr)} is primitive. The columns of
$n$ indexed by $\mathcal C$ are \m{\cc{S}{c}\tilde u\binom{z_1}{e}}; they are integral, and
\m{\widetilde{\lmat{\gamma}u}\binom{z_1}{e}=\tilde v\,e}. With Step~1 this shows \m{\lifts{S}{c}{\gamma}{\tilde v\,e}}.

\emph{Step 3: Jacobi's identity.} Let $F$ and $p'$ be the images under \m{\cc{S}{c}} of the first columns
($e_a$ for \m{a<\varrho}; $\gamma e_a$ for \m{a<j}, \m{a<\varrho}) and of the last columns ($e^*_a$ for
\m{a\in B}; $z$; $\gamma z$) of the matrix $P$ of Lemma~\ref{l:bordered}; by the last sentence of
§\ref{s:setting}(e), $F$ consists of the vectors $e_a$ and $ke_a$. The matrix $F$ is integral because \m{k\in\KS{S}}; \m{F^*J_fn=0} by
Step~1 and Input~1(d); and the sizes are complementary because \m{\#\mathcal C=\kk{\fs}+1} (Input~1(a)).
Jacobi's identity (Input~3(b)), for \m{J=J_f}, $F$, $n$, $p'$ and \m{c_0=N}, gives
\[\bigl[\det\bigl(n^*(J_f+N^{-1}J_fp'p'^*J_f)n\bigr)\bigr]\cdot\mathfrak i(F)^2\cdot[N]^{\#B+2}=[\det\tilde M]
=\bigl(N^{2\ninf{E}}\mnorm{S}{j}{\gamma}\bigr)^2,\]
the Gram matrix in the coordinates of the lattice being the rational one (Input~1(d)) and the last equality
being Lemma~\ref{l:bordered}. On the left, with $Q$ the matrix of the products of the columns of $u$ with
$e^*_a$ (\m{a\in B}), $z$ and~$\gamma z$,
\[n^*(J_f+N^{-1}J_fp'p'^*J_f)n=\bigl(\begin{smallmatrix}a&z_1\\0&e\end{smallmatrix}\bigr)^*\,
\widetilde{\bigl(u^*\Jrat{S}u+N^{-1}Q^*Q\bigr)}\,\bigl(\begin{smallmatrix}a&z_1\\0&e\end{smallmatrix}\bigr),\]
\[u^*\Jrat{S}u+N^{-1}Q^*Q=\begin{pmatrix}N^{-1}\cdot1_B&0\\0&v^*\lmet{S}v\end{pmatrix}.\]
The second identity follows from Step~1. A column of $u$ indexed by $B$ lies in \m{X\cap\gamma X}: its
products with $z$, $\gamma z$ and with the other columns vanish, and its products with the $e^*_a$,
\m{a\in B}, are its $e_a$-coordinates; this gives the block \m{N^{-1}\cdot1_B} and the zero blocks. For the
columns indexed by~$\mathcal C$, the products with the $e^*_a$, \m{a\in B}, vanish, the products with $z$ and
$\gamma z$ are the values of \m{\zp{T}} on the link images, which gives the term \m{N^{-1}\zp{T}^*\zp{T}} of the
link metric, and Input~1(b), with \m{\gamma'=\gamma} and the coefficients \m{\tfrac12}, \m{-\tfrac12}, gives
\m{u^*\Jrat{S}u=v^*\,\tfrac12\diag(T,T)\,v} on them. Taking generalised indices,
\[[\det a]^2\,[\det e]^2\cdot[N]^{-\#B}\cdot|N_{E/\Q}\det(v^*\lmet{S}v)|\cdot\mathfrak i(F)^2\cdot[N]^{\#B+2}
=[N]^2\,\mnorm{S}{j}{\gamma}^2,\]
that is
\begin{equation}\label{e:jacobi}
\ldsq{S}{v}{e}\cdot[\det a]^2\cdot\mathfrak i(F)^2=\mnorm{S}{j}{\gamma}^2 .
\end{equation}
\emph{Step 4: the indices.} The formal text shows two divisibilities. First, the matrix $a$ is integral, the
index of its column lattice in the integral vectors is \m{[\det a]\ne0}, and by the relation of Input~1(h),
applied to the columns of $n$ indexed by~$B$, and the description of the frame Gram matrix of $k$
(Input~1(d)), the lattice of the conjugated first $\varrho$ rows of $Z(k)$ is generated by that of its first
\m{\min(j,\varrho)} rows and the image of the column lattice of~$a$; so \m{\xri{S}{j}{k}} divides \m{[\det a]}.
Second, \m{\xri{S}{j}{k^{-1}}} divides the index of the column lattice of \m{\ifp{k}} in that of~$F$,
\m{\sat{S}{k}} is the index of the column lattice of \m{\ifp{k}} in the lattice of the integral vectors of its
span, and by the first assertion of Jacobi's identity that lattice lies in the lattice of the $y$ with $Fy$
integral; as indices multiply, \m{\xri{S}{j}{k^{-1}}\cdot\sat{S}{k}} divides \m{\mathfrak i(F)}, which is non-zero by
\eqref{e:jacobi} and \m{\mnorm{S}{j}{\gamma}>0}. Hence \m{w_0\ne0}, \m{w_0\le[\det a]\,\mathfrak i(F)}, and
\eqref{e:jacobi} gives \m{\ldsq{S}{v}{e}\,w_0^2\le\mnorm{S}{j}{\gamma}^2}.
\end{proof}
\leanline{\lean{existsliftslinkDiscSqle}}

\begin{corollary}[positivity]\label{c:pos}
In Setting~\ref{s:S}, suppose that $E$ is stably unramified. Let $c$, $c'$ be classes of~$S$, \m{s\in\GRes{S}},
$j$ a level and \m{\gamma\in\terms{S}{c}{s}{c'}} with \m{\haslevel{S}{j}{\gamma}}, and \m{k=\pos{S}{c}{c'}{\gamma}}.
Then \m{1\le\xri{S}{j}{k}}, \m{1\le\sat{S}{k}} and \m{0<\cgauge{S}{j}{c}{c'}{\gamma}}.
\end{corollary}
\begin{proof}
Apply Proposition~\ref{p:lifts} to a basis of the link: \m{w_0\ne0}, so its three factors, which are natural
numbers, are at least~$1$; and the numerator of the gauge is at least
\m{N^{\ninf{E}}\mnorm{S}{j}{\gamma}>0} (Input~1(e)).
\end{proof}
\leanline{\lean{onelexRowsIndexandsaturationIndex}}

\begin{lemma}[an integral matrix from lifts]\label{l:int}
In Setting~\ref{s:S}, let $c$ be a class of~$S$, \m{\gamma\in\GS{S}{E}}, $\mathcal C$ an index set, $m$ a matrix
over \m{\AEf} with rows \m{\Mx{\fs}\sqcup\Mx{\fs}} and columns~$\mathcal C$, and $c'$ a class of~$S$. Suppose that
\m{\lifts{S}{c}{\gamma}{m}}. Then, with \m{k=\pos{S}{c}{c'}{\gamma}}, there is an integral matrix $n$ with rows
\m{\Ix{\fs}} and columns $\mathcal C$ such that
\[\zp{\Tlat{S}{\AEf}}\,\lmat{k}\,n=\widetilde{\zp{T}}\,m,\qquad\ifp{k}^*J_fn=0,\]
and such that, for every matrix $p$ with rows~$\kappa$: if the first half of $m$ is $p$ on the rows $\kappa$
and $0$ on the row~$z$, then the middle rows of $n$ are \m{\ccX{S}{c}\,p}.
\end{lemma}
\begin{proof}
Let \m{(\mathcal C',t,w)} be as in the definition of integral lifts and put \m{n=\cc{S}{c}\tilde tw}. It is
integral. As \m{\cc{S}{c}\widetilde{\fpair{\gamma}}=\fpair{k}}, Input~1(d) turns the pairings of $n$ with the
columns of \m{\fpair{k}} into those of \m{\tilde tw} with the columns of \m{\fpair{\gamma}}: those with the $e_a$,
$\gamma e_a$ vanish, and those with $z$, $\gamma z$, which are the values of \m{\zp{\cdot}} on the link images,
are \m{\widetilde{\zp{T}}\,\widetilde{\lmat{\gamma}t}\,w=\widetilde{\zp{T}}\,m}. The last assertion is read off
from the block form of \m{\cc{S}{c}} on the middle coordinates.
\end{proof}
\leanline{\lean{existsisIntMatrix}}

\begin{lemma}[orthogonal splitting of lifts]\label{l:split}
In Setting~\ref{s:S}, suppose that $E$ is stably unramified. Let $c$ be a class of~$S$, \m{\gamma\in\GS{S}{E}},
$\mathcal C$ a finite index set, $m$ a matrix over $E$ with rows \m{\Mx{\fs}\sqcup\Mx{\fs}} and
columns~$\mathcal C$, $d$ a natural number, $y$ a matrix over $E$ with the same rows and $d$ columns, and $f$ an
invertible matrix over \m{\AEf} with rows and columns indexed by $\mathcal C$ and the $d$ columns of~$y$.
Suppose that \m{m^*\lmet{S}\,y=0} and that \m{\lifts{S}{c}{\gamma}{\widetilde{(m\,|\,y)}f}}. Then there are
\m{a\in\GL_{\mathcal C}(\AEf)}, \m{v\in\GL_d(\AEf)} and a matrix $z_1$ over \m{\AEf} with rows $\mathcal C$ and $d$
columns such that
\[\lifts{S}{c}{\gamma}{\tilde ma},\qquad\lifts{S}{c}{\gamma}{\tilde mz_1+\tilde yv},\qquad
\ldsq{S}{m}{a}\cdot\ldsq{S}{y}{v}=\ldsq{S}{(m\,|\,y)}{f} .\]
\end{lemma}
\begin{proof}
By Input~3(a), applied to~$f^{-1}$, there are $a$, $v$, $z_1$ such that
\m{f^{-1}\bigl(\begin{smallmatrix}a&z_1\\0&v\end{smallmatrix}\bigr)} is primitive, in particular integral with
\m{[\det a]\,[\det v]=[\det f]}. Let \m{(\mathcal C',t,w)} be lift data for \m{\widetilde{(m\,|\,y)}f}. Then
\m{(\tilde ma\,|\,\tilde mz_1+\tilde yv)=\widetilde{\lmat{\gamma}t}\cdot w\,f^{-1}\bigl(\begin{smallmatrix}a&z_1\\0&v\end{smallmatrix}\bigr)},
and \m{\cc{S}{c}\tilde t\cdot w\,f^{-1}\bigl(\begin{smallmatrix}a&z_1\\0&v\end{smallmatrix}\bigr)} is integral; taking
the two blocks of columns gives the two lifts. Since \m{m^*\lmet{S}y=0}, the Gram matrix of \m{(m\,|\,y)} for
the link metric is block diagonal, so its determinant is the product of those of $m$ and of~$y$; with
\m{[\det f]^2=[\det a]^2[\det v]^2} this is the identity of the discriminants.
\end{proof}
\leanline{\lean{existsorthogonalsplit}}

\subhead{The vector that carries the term discriminant.}
\begin{lemma}\label{l:vector}
In Setting~\ref{s:S}, let $j$ be a level and \m{\gamma\in\GS{S}{E}} with \m{\haslevel{S}{j}{\gamma}}. Then there
are \m{a\in E^\kappa} and \m{y\in E} such that the pair \m{(a-z,\ yz)} of middle vectors lies in
\m{\link{\gamma}}, and
\[\tdisc{S}{\gamma}=\lminor{\Zf{S}{\gamma}}{j}\cdot\overline{\lminor{\Zf{S}{\gamma}}{j}}\cdot(1-\bar yy).\]
\end{lemma}
\begin{proof}
It suffices to find a coefficient vector $x$ such that \m{t=\sum_bx_b\gamma\mathfrak f_b} has middle part
\m{a-z} and such that \m{\gamma^{-1}t=\sum_bx_b\mathfrak f_b} has middle part $yz$, that is \m{y=x_\varrho}, with
$t$ orthogonal to $X$ and~$\gamma X$: then the pair is the link image of~$t$.

If \m{j=\varrho+1}, then $D=\det\Zf{S}{\gamma}=\lminor{\Zf{S}{\gamma}}{j}\ne0$; let $x'$ be the solution of
\m{\Zf{S}{\gamma}x'=D\,e_\varrho} and \m{x=(N/D)\,x'}, so that \m{\Zf{S}{\gamma}x=N\,e_\varrho}. By Input~1(f) the
vector \m{t=\sum_bx_b\gamma\mathfrak f_b} has no coordinates at the $e^*_a$ and has \m{t_z=-1}; so its middle part
is \m{a-z} with \m{a\in E^\kappa}, and it is orthogonal to~$X$. It is orthogonal to $\gamma X$ because
\m{h(\gamma e_a,t)=h(e_a,\sum_bx_b\mathfrak f_b)=0}. By Cramer's rule \m{x'_\varrho=\lminor{\Zf{S}{\gamma}}{\varrho}},
the cofactor of the top entry; so \m{y=x_\varrho=N\lminor{\Zf{S}{\gamma}}{\varrho}/D}, and
\[\tdisc{S}{\gamma}=D\bar D-N^2|\lminor{\Zf{S}{\gamma}}{\varrho}|^2=D\bar D(1-\bar yy).\]

If \m{j\le\varrho}, the normal form of \m{\gamma^{-1}} for the prescription $-1$ at the top (Input~1(e)) gives
$x$, $y'$ with \m{y'_\varrho=-1} and \m{\sum_bx_b\mathfrak f_b=\gamma^{-1}\sum_by'_b\mathfrak f_b}; take
\m{t=\sum_by'_b\mathfrak f_b}, whose middle part is \m{-z} (so \m{a=0}), and \m{y=x_\varrho}, for which
\m{\bar yy=\bar y'_\varrho y'_\varrho=1}. Both $t$ and \m{\gamma^{-1}t} lie in the span of the frame, so $t$ is
orthogonal to $X$ and to~$\gamma X$. As \m{j\le\varrho}, \m{\tdisc{S}{\gamma}=0} by the definition of the level,
and \m{1-\bar yy=0}.
\end{proof}
\leanline{\lean{existsmemlinktermDisceq}}

\subhead{Comparison of the minor with a partial discriminant.}
\begin{proposition}\label{p:minor}
In Setting~\ref{s:S}, suppose that $E$ is stably unramified. Let $c$, $c'$ be classes of~$S$, \m{s\in\GRes{S}} an
element that agrees, $x$ a natural number such that $s$ fixes norm at least~$x$, $j$ a level and
\m{\gamma\in\terms{S}{c}{s}{c'}} with \m{\haslevel{S}{j}{\gamma}}. Let $r$ and $i$ be natural numbers, $w$ a
matrix over \m{\AEf} with rows \m{\Mx{\fs}\sqcup\Mx{\fs}} and $r$ columns, $y$ a matrix over $E$ with the same
rows and $r$ columns, and \m{v\in\GL_r(\AEf)}. Suppose that
\begin{enumerate}[label=\textup{(\alph*)},leftmargin=2.3em,itemsep=1pt]
\item if \m{i=0}, then \m{\dim(\link{\gamma}\cap\ker\zp{T})+r=\dim\link{\gamma}};
\item the columns of $y$ lie in \m{\link{\gamma}}, and \m{\zp{T}\,y} has a left inverse;
\item \m{\lifts{S}{c}{\gamma}{w}}, and \m{\widetilde{\zp{T}}\,w=\widetilde{\zp{T}y}\,v}.
\end{enumerate}
Then \m{r\le1+(j-\rh{\fs})_+} and, with \m{x'=x} if \m{\rh{\fs}<j} and \m{i=0}, and \m{x'=1} otherwise,
\[x'^2\cdot\mnorm{S}{j}{\gamma}^2\ \le\ \ldsq{S}{y}{v}\cdot\max\bigl(\tdnorm{S}{\gamma},\ \NS{S}^{\ninf{E}}\mnorm{S}{j}{\gamma}\bigr)^2 .\]
\end{proposition}
\begin{proof}
Let \m{k=\pos{S}{c}{c'}{\gamma}\in\KS{S}} and let $n$ be the integral matrix of Lemma~\ref{l:int} for~$w$.
By Lemma~\ref{l:vector} for $\gamma$ and for \m{\gamma^{-1}} there are pairs \m{\varepsilon_1=(a-z,\ uz)} in
\m{\link{\gamma}} with \m{\tdisc{S}{\gamma}=D_j\bar D_j(1-\bar uu)}, \m{D_j=\lminor{\Zf{S}{\gamma}}{j}}, and
\m{(a'-z,\ u'z)} in \m{\link{\gamma^{-1}}}; exchanging the two components of the latter gives a second pair
\m{\varepsilon_2=(u'z,\ a'-z)} in \m{\link{\gamma}}. Let $\varepsilon$ be the matrix with the columns
$\varepsilon_1$, $\varepsilon_2$.

\emph{Two identities.} For every matrix $w'$ with columns in \m{\link{\gamma}},
\begin{equation}\label{e:eps}
\varepsilon^*\,\lmet{S}\,w'=\zp{T}\,w' ,\qquad\zp{T}\,\varepsilon=N\begin{pmatrix}1&-\bar u\\-u&1\end{pmatrix}.
\end{equation}
Indeed the value of \m{\zp{T}} on a pair with coordinates $s_1$, $s_2$ at $z$ is \m{(-Ns_1,-Ns_2)}, which gives
$\zp{T}\varepsilon$ with the entry \m{-Nu'} in the upper right corner; and $\varepsilon_1^*\lform{S}\varepsilon_2=0$
(Input~1(a)) reads \m{Nu'-N\bar u=0}, so \m{u'=\bar u}. For the first identity write
$\lmet{S}=\tfrac12\lform{S}+\diag(0,T)+N^{-1}\zp{T}^*\zp{T}$; the link form vanishes on the link, and the two
other terms, computed from the coordinates of $\varepsilon_1$, $\varepsilon_2$ and using once more that
$\varepsilon_2^*\lform{S}w'=0$, add up to \m{\zp{T}w'}.

\emph{A minor.} As \m{\zp{T}y} has two rows and a left inverse, \m{r\le2}, and there is an injective map
$\varsigma_0$ from the $r$ columns to the two rows such that the corresponding minor of \m{\zp{T}y} is non-zero.
Let $\varepsilon_{\varsigma_0}$ be the matrix of the columns of $\varepsilon$ selected by~$\varsigma_0$; by
\eqref{e:eps} the matrix \m{\varepsilon_{\varsigma_0}^*\lmet{S}y} is that minor's matrix, so its determinant is
non-zero.

\emph{Liouville's inequality} (Input~2, with \m{Q=\lmet{S}}, which is positive definite at the infinite
places, \m{a=\varepsilon_{\varsigma_0}}, \m{a'=y}, the identity and~$v$) gives
\begin{equation}\label{e:liou}
x'^2\ \le\ |N_{E/\Q}\det(\varepsilon_{\varsigma_0}^*\lmet{S}\varepsilon_{\varsigma_0})|\cdot\ldsq{S}{y}{v},
\end{equation}
provided \m{x'\le[\delta]} whenever \m{\delta=\det\bigl((\widetilde{\zp{T}y}\,v)_{\varsigma_0}\bigr)=
\det\bigl((\widetilde{\zp{T}}w)_{\varsigma_0}\bigr)} is a unit. By Lemma~\ref{l:int},
\m{\widetilde{\zp{T}}w} is the matrix of the pairings of the columns of the integral matrix $n$ with $z$
and~$kz$, and the pairings of $n$ with the $e_a$, $ke_a$ vanish; these are the pairings $w^*J_fv$ of
Input~4. If \m{\varrho<j} and \m{i=0}, then
\m{\tdisc{S}{\gamma}\ne0}, so \m{1-\bar uu\ne0} and \m{\zp{T}\varepsilon} is invertible; hence
\m{\dim\link{\gamma}=2+\dim(\link{\gamma}\cap\ker\zp{T})}, hypothesis (a) gives \m{r=2}, the minor is the full
determinant, and Input~4 gives \m{x\le[\delta]}; otherwise \m{x'=1\le[\delta]} because
$\delta$ is a unit that is the determinant of an integral matrix.

\emph{The determinant of the selected block.} By \eqref{e:eps},
$\varepsilon^*\lmet{S}\varepsilon=\zp{T}\varepsilon$. If \m{r=0}, the determinant in \eqref{e:liou} is $1$. If
\m{r=1}, it is a diagonal entry, equal to~$N$, of absolute norm \m{N^{2\ninf{E}}}. If \m{r=2}, it is
\m{N^2(1-\bar uu)}, and it is non-zero: $\zp{T}y=\varepsilon^*\lmet{S}y$ is invertible, so $\varepsilon$ has
rank~$2$ and the link metric is positive definite. Hence $\tdisc{S}{\gamma}=D_j\bar D_j(1-\bar uu)\ne0$ and
\m{j=\varrho+1} by the definition of the level; so $r\le1+(j-\varrho)_+$ in all cases. Moreover, as
$|N_{E/\Q}(\xi)|=\prod_\tau|\xi_\tau|^2$ and $\mnorm{S}{\varrho+1}{\gamma}=|N_{E/\Q}(D_j)|\,N^{-2\ninf{E}}$,
$|N_{E/\Q}(N^2(1-\bar uu))|\cdot\mnorm{S}{j}{\gamma}^2=|N_{E/\Q}(D_j)|^2\,|N_{E/\Q}(1-\bar uu)|=\tdnorm{S}{\gamma}^2$. Multiplying \eqref{e:liou} by $\mnorm{S}{j}{\gamma}^2$: for
\m{r\le1} the right side is at most $\ldsq{S}{y}{v}\,(N^{\ninf{E}}\mnorm{S}{j}{\gamma})^2$, and for \m{r=2} it
is $\ldsq{S}{y}{v}\,\tdnorm{S}{\gamma}^2$. Both are bounded by the maximum.
\end{proof}
\leanline{\lean{minorNormsqlelinkDiscSqmul}}

\subhead{Proof of Theorem B.}
\begin{proof}[Proof of Theorem B]
Write \m{k_0=\kk{\fs}}, $k_\gamma=\pos{S}{c}{c'}{\gamma}$, $L_0=L\cap\ker\zp{T}$ and
$T^\circ=\clink{\gamma}{\gamma'}$.

\emph{Bases.} Let $m$ be a matrix whose columns, indexed by a finite set~$\mathcal C$, are a basis of~$L_0$,
and $p$ its part on the rows $\kappa$ of the first components. The rows $z$ of $m$ vanish, and the link form
vanishes on~$L_0$; so \m{m^*\lmet{S}m=p^*p}, and \m{T^\circ} is the column span of~$p$, of dimension
\m{\#\mathcal C}. By Input~3(c) (with \m{A=m^*\lmet{S}}, whose product with $m$ is invertible by Input~1(c)),
extend $m$ by a matrix $u$ to a basis \m{(m\,|\,u)} of~$L$ with \m{m^*\lmet{S}u=0}, and then \m{(m\,|\,u)} by a
matrix $y$ with $i$ columns to a basis of \m{\link{\gamma}} with \m{(m\,|\,u)^*\lmet{S}y=0} and
\m{y\,t\in L} only for \m{t=0}. Then \m{i=k_0+1-\dim L} by Input~1(a), and if \m{i=0} then \m{L=\link{\gamma}}.
Moreover \m{\zp{T}u} has a left inverse, the columns of $u$ being independent modulo~$L_0$.

\emph{Lifts.} By Proposition~\ref{p:lifts} there is $e$ with
$\lifts{S}{c}{\gamma}{\widetilde{(m\,|\,u\,|\,y)}e}$, $\ldsq{S}{(m\,|\,u\,|\,y)}{e}\,w_0^2\le\mnorm{S}{j}{\gamma}^2$
and $w_0=\xri{S}{j}{k_\gamma}\xri{S}{j}{k_\gamma^{-1}}\sat{S}{k_\gamma}\ne0$. Lemma~\ref{l:split}, applied to \m{m_1=(m\,|\,u)} and~$y$,
gives $a'$, $v$, $z$ with integral lifts of \m{\tilde m_1a'} and of $\tilde m_1z+\tilde yv$; applied to $m$
and~$u$ and the lift \m{\tilde m_1a'}, it gives $a$, $g_0$, $z'$ with integral lifts of \m{\tilde ma} and of
$\tilde mz'+\tilde ug_0$; and the discriminants multiply:
\begin{equation}\label{e:three}
\ldsq{S}{m}{a}\cdot\ldsq{S}{u}{g_0}\cdot\ldsq{S}{y}{v}=\ldsq{S}{(m\,|\,u\,|\,y)}{e}.
\end{equation}
The matrix of the statement of Theorem~B is $m_1$, whose columns span~$L$, with the index set
\m{\mathcal C\sqcup\{1,\dots,r\}}, $r$ the number of columns of~$u$.

\emph{The part in~$L_0$.} Let $n$ be the integral matrix of Lemma~\ref{l:int} for \m{\tilde ma} and~$c'$. Its
middle rows are \m{\ccX{S}{c}\,\tilde p\,a}, its link image is killed by \m{\zp{\cdot}}, and its columns are
orthogonal to \m{X+k_\gamma X}; so by Input~1(h) the values \m{(h(k_\gamma\mathfrak f_l,\cdot))_{l<\varrho}} of its middle
part lie in the lattice of the conjugated first $\varrho$ rows of \m{Z(k_\gamma)}. Input~3(d), with \m{V=T^\circ},
$p$, $a$ and \m{D=[\det(p^*p)]\,[\det a]^2=\ldsq{S}{m}{a}\ne0}, gives~(i) and
\begin{equation}\label{e:disc}
\bigl(\ldisc{S}{c}{T^\circ}\cdot\fvi{S}{c}{k_\gamma}{1}{T^\circ}\bigr)^2\ \le\ \ldsq{S}{m}{a}.
\end{equation}
\emph{The part in $L$ off $L_0$.} Proposition~\ref{p:minor}, for the matrix $u$, the
lift \m{\tilde mz'+\tilde ug_0} (whose image under \m{\widetilde{\zp{T}}} is that of \m{\tilde ug_0}) and the
number~$i$ (if \m{i=0} then \m{L=\link{\gamma}} and hypothesis (a) holds), gives \m{r\le1+(j-\varrho)_+} and
\begin{equation}\label{e:G}
x'^2\,\mnorm{S}{j}{\gamma}^2\ \le\ \ldsq{S}{u}{g_0}\cdot\max\bigl(\tdnorm{S}{\gamma},N^{\ninf{E}}\mnorm{S}{j}{\gamma}\bigr)^2 .
\end{equation}
Since \m{k_0+1=\dim\link{\gamma}=\#\mathcal C+r+i}, we get \m{k_0\le\dim T^\circ+i+(j-\varrho)_+}.

\emph{Assembly.} Multiplying \eqref{e:disc}, \eqref{e:G} and \m{\ldsq{S}{y}{v}}, then using \eqref{e:three} and
the bound for the full lift,
\[\begin{aligned}\bigl(\ldisc{S}{c}{T^\circ}\fvi{S}{c}{k_\gamma}{1}{T^\circ}\bigr)^2x'^2\,\mnorm{S}{j}{\gamma}^2\,\ldsq{S}{y}{v}
&\le\max(\cdots)^2\,\ldsq{S}{(m\,|\,u\,|\,y)}{e}\\&\le\max(\cdots)^2\,\mnorm{S}{j}{\gamma}^2\,w_0^{-2}.\end{aligned}\]
Dividing by $\mnorm{S}{j}{\gamma}^2>0$ and multiplying by $\sat{S}{k_\gamma}^2$, the left side becomes
\[\bigl(\pinv{S}{c}{c'}{\gamma}{\gamma'}\,x'\bigr)^2\,\ldsq{S}{y}{v},\]
and the right side becomes
\[\max(\cdots)^2\cdot\bigl(\xri{S}{j}{k_\gamma}\,\xri{S}{j}{k_\gamma^{-1}}\bigr)^{-2}=\cgauge{S}{j}{c}{c'}{\gamma}^2 .\qedhere\]
\end{proof}

\subhead{Proof of Theorem A.}
\begin{proof}[Proof of Theorem A]
Apply Theorem~B to \m{(\gamma,\gamma')} and to \m{(\gamma',\gamma)}; the space $L$ is the same, so both give the
same number \m{i=\kk{\fs}+1-\dim L} and the same~$x'$. This gives (i) and (ii) of Theorem~A, matrices $m$, $y$,
$v$, $z$ for $\gamma$ and $m'$, $y'$, $v'$, $z'$ for~$\gamma'$, with $i$ columns in $y$ and in~$y'$, and
\begin{equation}\label{e:two}
\begin{aligned}\bigl(\pinv{S}{c}{c'}{\gamma}{\gamma'}x'\bigr)^2\ldsq{S}{y}{v}&\le\cgauge{S}{j}{c}{c'}{\gamma}^2,\\
\bigl(\pinv{S}{c}{c'}{\gamma'}{\gamma}x'\bigr)^2\ldsq{S}{y'}{v'}&\le\cgauge{S}{j}{c}{c'}{\gamma'}^2 .\end{aligned}
\end{equation}
\emph{The pairing of the complements.} Let $Q_0$ be the diagonal matrix with \m{\lmet{S}Q_0=\lform{S}}; it
satisfies \m{Q_0^*\lform{S}=4\,\lmet{S}} (both are checked on the diagonal entries). We apply Input~2 with
\m{Q=\lmet{S}}, \m{a=\gS{S}^{-1}y}, \m{a'=Q_0y'}, $v$, $v'$ and \m{x_0=1}. Here
\[a^*Qa'=\gS{S}^{-1}y^*\lform{S}y',\qquad a^*Qa=\gS{S}^{-2}y^*\lmet{S}y,\qquad a'^*Qa'=4\,y'^*\lmet{S}y',\]
and the determinants of the last two have the absolute norms
\[\gS{S}^{-4i\ninf{E}}\,|N_{E/\Q}\det(y^*\lmet{S}y)|\qquad\text{and}\qquad2^{4i\ninf{E}}\,|N_{E/\Q}\det(y'^*\lmet{S}y')| .\]
So, once its hypotheses are verified, Input~2 gives
\begin{equation}\label{e:pair}
1\ \le\ \Bigl(\frac{2}{\gS{S}}\Bigr)^{4i\ninf{E}}\ldsq{S}{y}{v}\cdot\ldsq{S}{y'}{v'} .
\end{equation}
\emph{Invertibility.} Suppose \m{y^*\lform{S}y'\,t=0}. The vector \m{y't} lies in \m{\link{\gamma'}}, which is
isotropic for the link form and contains~$L$; so \m{y't} is orthogonal to the columns of $m$ and of~$y$, which
span \m{\link{\gamma}}. As the link is Lagrangian (Input~1(a)), \m{y't\in\link{\gamma}}, hence \m{y't\in L} and
\m{t=0} by the property of~$y'$.

\emph{Integrality.} Let \m{(t_1,w_1)} and \m{(t_2,w_2)} be lift data for \m{\tilde mz+\tilde yv} and for
\m{\tilde m'z'+\tilde y'v'}. The columns of $m$, $m'$ lie in~$L$, which is orthogonal for the link form to
\m{\link{\gamma}} and to \m{\link{\gamma'}}; so
\[v^*\,\widetilde{y^*\lform{S}y'}\,v'=(\tilde mz+\tilde yv)^*\,\widetilde{\lform{S}}\,(\tilde m'z'+\tilde y'v')
=w_1^*\,\widetilde{t_1^*\Jrat{S}(1-\gamma\gamma'^{-1})t_2}\,w_2,\]
the last step by Input~1(b) with the coefficients $1$,~$1$. In the coordinates of the lattice of $c$ this is
\m{(\cc{S}{c}\tilde t_1w_1)^*J_f(1-\kappa_0)(\cc{S}{c}\tilde t_2w_2)} with
\m{\kappa_0=\pos{S}{c}{c}{\gamma\gamma'^{-1}}}; as \m{\gamma\gamma'^{-1}} is a term over \m{(c,1,c)},
\m{\gS{S}^{-1}(1-\kappa_0)} is integral (Input~1(g)), and the other factors are integral. Hence the matrix
\m{v^*\,\widetilde{a^*Qa'}\,v'} is integral, and if its determinant $\delta$ is a unit then \m{1\le[\delta]}.

\emph{Conclusion.} Multiplying the two inequalities \eqref{e:two} and using \eqref{e:pair},
\[\bigl(\pinv{S}{c}{c'}{\gamma}{\gamma'}\pinv{S}{c}{c'}{\gamma'}{\gamma}\,x'^2\bigr)^2\le
\Bigl(\frac{2}{\gS{S}}\Bigr)^{4i\ninf{E}}\bigl(\cgauge{S}{j}{c}{c'}{\gamma}\cgauge{S}{j}{c}{c'}{\gamma'}\bigr)^2,\]
and all the quantities are non-negative; taking square roots gives~(iii).
\end{proof}

\sect{Where the results are used}{s:used}
Theorems A and~B (the latter for its part~(i)) and Corollary~\ref{c:pos} are used in the crude estimate of
the note on the energy of amplitude sums of this unit, in the estimate of the pairs of terms of one level
(Input~3 of Note~F); Corollary~\ref{c:pos} is used also in the modules on thin classes and on the block bound.
Proposition~\ref{p:lifts}, Lemmas \ref{l:bordered}, \ref{l:int}, \ref{l:split}, \ref{l:vector} and
Proposition~\ref{p:minor} are used only inside this module. Nothing of this note is used directly by the
filling theorem.

\sect{Notes on the formal text}{s:notes}
\begin{enumerate}[leftmargin=1.6em,itemsep=3pt]
\item \emph{Arguments that the census reads differently from Lean; to be settled by} \texttt{\#check}
\emph{on a build.} In Lemma~\ref{l:int} the formal statement names the index set of the columns of~$m$, a
variable of the file, so Lean takes it as an argument although the census does not list it; it is the set
$\mathcal C$ of the statement here. For Theorem~B the census lists a variable of the file for a matrix, its
index set and the finiteness of that set; the statement introduces its own matrix $m$ after ``there are'', and
Lean takes none of the three. They have no counterpart in the statement of Theorem~B.
\item \emph{Statements cut by the census.} The statements of Lemma~\ref{l:bordered},
Proposition~\ref{p:lifts}, Lemma~\ref{l:int} and Theorem~B contain local definitions at which the census ends
the statement; they were read to their true end, with the second reader's table of corrected spans.
\item \emph{The columns of the bordered matrix.} In Lemma~\ref{l:bordered} the columns are $\gamma e_a$ for
\m{a<j} and $e^*_a$ for \m{a\ge j}, as the code selects them; the comment of the source says \m{a\le j} and
\m{a>j}. The second reader of the dossier had noted this.
\item \emph{The matrix $m$ of Theorem~B} spans \m{L=\link{\gamma}\cap\link{\gamma'}}, not its part in
\m{\ker\zp{T}}; the latter appears only in the proof.
\item \emph{Truncated differences.} \m{(j-\rh{\fs})_+} and \m{\kk{\fs}+1-\dim L} are differences of natural
numbers in the formal text; the first is $0$ or~$1$, and the second is not truncated since
\m{\dim L\le\kk{\fs}+1}.
\item \emph{Exponents computed here.} In Proposition~\ref{p:lifts} the powers of $N$ in Jacobi's identity
cancel exactly: \m{[N]^{\#B+2}\cdot[N]^{-\#B}=[N]^2=N^{4\ninf{E}}}, the square of the factor in
Lemma~\ref{l:bordered}. In Theorem~A the factor \m{(2/\gS{S})^{2i\ninf{E}}} is the square root of
\m{\gS{S}^{-4i\ninf{E}}\cdot2^{4i\ninf{E}}}, the absolute norms of \m{\gS{S}^{-2i}} and of~$4^{i}$. In
Proposition~\ref{p:minor} the case \m{r=2} uses \m{|N_{E/\Q}(N^2)|=N^{4\ninf{E}}}, which cancels the
normalisation of \m{\mnorm{S}{\varrho+1}{\gamma}}.
\item \emph{What was not re-derived.} One passage is stated as the formal text has it and was not derived
again for this note: the two divisibilities of Step~4 of Proposition~\ref{p:lifts}. All other steps are proved
in the text from the inputs. In Lemma~\ref{l:bordered} the formal text reorders the rows and columns of the
bordered matrix and verifies its block form entry by entry; the note gives the three lines that this
verification amounts to.
\item \emph{The threshold $x$} of ``fixes norm at least~$x$'' is a natural number in the formal text; in the
inequalities of Theorems A and~B and of Proposition~\ref{p:minor} it is read as a real number.
\item \emph{Input~4.} The formal statement is about an integral matrix and its plain products with the frame
and its image; the note states it for the pairings of the hermitian form, as it is applied.
\end{enumerate}

\subhead{How this note was written.} From the structured proofs of the module on link pairings in the fourth
chunk of the unit's dossier, which a second reader had re-read completely. I read the module in the Lean text
line by line and compared every statement of this note with it argument by argument; I read in the Lean text
the definitions recalled in §\ref{s:setting}. The inputs are stated from the dossier's readings of their quoted
statements; of them I read in the Lean text only the definitions they involve. The block form of the bordered matrix in Lemma~\ref{l:bordered} is
written as the referee of this note gave it. Arguments that are mine and
replace matrix computations of the formal text: the evaluation of \m{\det A} in Lemma~\ref{l:bordered} by a
rank-one update, the block identity in Step~3 of Proposition~\ref{p:lifts}, the two identities \eqref{e:eps},
the exclusion of \m{r=2} for \m{j\le\varrho}, and the computation of the pairing in Theorem~A. Nothing was
compiled for this note. The program check of the unit's dossier covered \UnitSteps{} leaf steps of \UnitProofs{}
proofs.

\head{Appendix A. Concordance with the Lean artifact}
The formal names of the objects and the letters of this note: \texttt{link}, \texttt{commonLink} and
\texttt{linkMatrix} are \m{\link{\gamma}}, \m{\clink{\gamma}{\gamma'}} and \m{\lmat{\gamma}}; \texttt{framePair} and
\texttt{isotropicFramePair} are \m{\fpair{\gamma}} and \m{\ifp{\gamma}}; \texttt{linkForm},
\texttt{linkMetric}, \texttt{zFormPair} are \m{\lform{S}}, \m{\lmet{S}}, \m{\zp{T}}; \texttt{linkDiscSq},
\texttt{minorNorm}, \texttt{classGauge}, \texttt{pairInvariant}, \texttt{termDiscNorm} are \m{\ldsq{S}{v}{e}},
\m{\mnorm{S}{j}{\gamma}}, \m{\cgauge{S}{j}{c}{c'}{\gamma}}, \m{\pinv{S}{c}{c'}{\gamma}{\gamma'}}, \m{\tdnorm{S}{\gamma}};
\texttt{xRowsIndex}, \texttt{saturationIndex}, \texttt{latticeDisc}, \texttt{frameValueIndex} are
\m{\xri{S}{j}{k}}, \m{\sat{S}{k}}, \m{\ldisc{S}{c}{V}}, \m{\fvi{S}{c}{k}{a}{V}}; \texttt{HasIntegralLifts},
\texttt{HasLevel} are \m{\lifts{S}{c}{\gamma}{m}}, \m{\haslevel{S}{j}{\gamma}}; \texttt{classCoord},
\texttt{classCoordX₀} are \m{\cc{S}{c}}, \m{\ccX{S}{c}}; \texttt{genIndex} is \m{[\xi]}. Two statements as the
artifact has them:
\quoted{pairInvariantmulle}
\quoted{existslinksplitleclassGauge}
{\fontsize{8.6}{11.6}\selectfont\renewcommand{\leanfont}{\ttfamily\fontsize{7.7}{10}\selectfont\color{leancol}}
\setlength{\LTleft}{0pt}\setlength{\LTright}{\fill}\setlength{\tabcolsep}{0pt}
\begin{longtable}{@{}p{27mm}>{\raggedright\arraybackslash}p{68mm}p{34mm}r@{}}
\textbf{In this note} & \textbf{Lean declaration} & \textbf{File} & \textbf{Lines}\\\midrule\endhead
\ConcordanceRows
\end{longtable}}

\head{Notation}
\begingroup\fontsize{8.6}{10.9}\selectfont\raggedright
\setlength{\columnsep}{9mm}\setlength{\parskip}{0pt}
\begin{multicols}{2}
\nt{$E$, $\tau$, $x_\tau$, \m{\ninf{E}}, \m{N_{E/\Q}}}{a CM field with its CM type, an infinite place, the image
at it, the number of infinite places, the norm}
\nt{\m{\AEf}, \m{\hO}, \m{\tilde a}, \m{[\xi]}}{finite adeles, integral adeles, the image of a matrix over~$E$,
the generalised index}
\nt{\m{\fs=(\varrho,\kappa,k_0)}}{a shape; \m{\rh{\fs}}, \m{\kk{\fs}}, \m{\Ix{\fs}}, \m{\Mx{\fs}}; $j$ a level; in
proofs \m{\varrho}, \m{k_0} abbreviate \m{\rh{\fs}}, \m{\kk{\fs}}}
\nt{$e_a$, $e^*_a$, $z$, \m{\mathfrak f_a}, $X$}{basis vectors, the frame, the span of the $e_a$; $h$ the
hermitian form}
\nt{$S$}{filling data: \m{\NS{S}} (also $N$), \m{\gS{S}}, \m{\GRes{S}}, \m{\Cl{S}}, \m{\GS{S}{E}}, \m{\GFin{S}},
\m{\KS{S}}, \m{\redS{S}}, \m{\KgS{S}}}
\nt{$T$, \m{\Trat{S}{E}}, \m{\Tlat{S}{\AEf}}, \m{\Jrat{S}}, $J_f$}{the middle form and the split Gram matrices,
rational and in the lattice basis}
\nt{$c$, $c'$, $s$, $t$, \m{\gamma}, \m{\gamma'}}{classes, elements of \m{\GRes{S}}, rational elements}
\nt{\m{\pos{S}{c}{c'}{\gamma}}, \m{\terms{S}{c}{s}{c'}}, $k$, \m{k_\gamma}, \m{\kappa_0}}{position, terms; positions
in proofs}
\nt{\m{\Zf{S}{\gamma}}, $Z(k)$, \m{Z_X(k)}, \m{\lminor{Z}{l}}, \m{D_j}, $D$}{frame Gram matrices, their first rows,
leading minors}
\nt{\m{\tdisc{S}{\gamma}}, \m{\tdnorm{S}{\gamma}}}{term discriminant and the product of its absolute values}
\nt{\m{\fpair{\gamma}}, \m{\ifp{\gamma}}, \m{\lmat{\gamma}}, \m{\link{\gamma}}, \m{\clink{\gamma}{\gamma'}}}{frame
pair, isotropic frame pair, link matrix, link, common link; the italic $L$, $L_0$ and \m{T^\circ} are subspaces
in the proof of Theorem~B (the upright \m{\mathrm L} belongs to the link matrix and to \m{\zp{T}})}
\nt{\m{\zp{T}}, \m{\lform{S}}, \m{\lmet{S}}, $Q_0$}{the $z$-pairings, link form, link metric, a diagonal matrix}
\nt{\m{\ldsq{S}{v}{e}}, \m{\mnorm{S}{j}{\gamma}}, \m{\cgauge{S}{j}{c}{c'}{\gamma}}, \m{\pinv{S}{c}{c'}{\gamma}{\gamma'}}}{squared
discriminant, minor norm, gauge, invariant of a pair}
\nt{\m{\cc{S}{c}}, \m{\ccX{S}{c}}, \m{\lifts{S}{c}{\gamma}{m}}}{class coordinates; integral lifts, for the
filling data~$S$ of the setting}
\nt{\m{\xri{S}{j}{k}}, \m{\sat{S}{k}}, \m{\ldisc{S}{c}{V}}, \m{\fvi{S}{c}{k}{a}{V}}, \m{\mathfrak i(F)}, $w_0$}{indices
and discriminants of lattices}
\nt{\m{\haslevel{S}{j}{\gamma}}, \m{\mathrm{pl}_j(Z)}, \m{\upsilon}}{level; the padded leading matrix; the scalar of a
normal form; \m{\Agr{S}{s}}, \m{\FixN{s}{x}} the predicates ``agrees'' and ``fixes norm at least $x$''}
\nt{$P$, $M$, $A$, $Z'$, $B$, $F$, $p'$, $Q$, $u$, $n$, $a$, $e$, $z_1$}{matrices of Lemma~\ref{l:bordered} and
Proposition~\ref{p:lifts}; here $P$ is a matrix, not the parameters}
\nt{$\mathcal C$, $\mathcal C'$, $m$, $y$, $v$, $z$, $f$, $d$, $r$, $i$, $w$, $x$, $x'$, $x_0$}{index sets, matrices
and numbers of the statements}
\nt{\m{\varepsilon}, \m{\varepsilon_1}, \m{\varepsilon_2}, $u'$, $a'$, \m{\varsigma_0}, \m{\delta}}{the two link
vectors, a selection of rows, a determinant, in Proposition~\ref{p:minor}}
\nt{$p$, $g_0$, $z'$, $m_1$, $m'$, $y'$, $v'$, $t_1$, $t_2$, $w_1$, $w_2$}{in the proofs of Theorems B and~A}
\nt{$J$, $c_0$, $d_p$, $V$, $q$, $l$, $b$, $D_0$, $O$, $G$, $K$, $U$, $R$, $H$, $I$, \m{\mathrm{fm}(k)}, $n_X$, $n_M$}{in
the inputs and as stated where they occur; \m{\GL}, \m{\Q}}
\end{multicols}
\endgroup

\end{document}
